\documentclass[11pt]{amsart}
\usepackage[T1]{fontenc}
\usepackage{lmodern,amsmath,amssymb,amsthm,mathtools,microtype,booktabs,array,longtable}
\usepackage[margin=1.06in]{geometry}
\usepackage[hidelinks]{hyperref}
\hypersetup{pdftitle={Input-Dependent Boundary Coding and the Full Error Range},pdfauthor={Qian Qi}}
\emergencystretch=2em
\numberwithin{equation}{section}
\newtheorem{theorem}{Theorem}[section]
\newtheorem{lemma}[theorem]{Lemma}
\newtheorem{proposition}[theorem]{Proposition}
\newtheorem{corollary}[theorem]{Corollary}
\theoremstyle{definition}\newtheorem{definition}[theorem]{Definition}
\theoremstyle{remark}\newtheorem{remark}[theorem]{Remark}
\newcommand{\Q}{\mathbb Q}\newcommand{\R}{\mathbb R}\newcommand{\C}{\mathbb C}\newcommand{\N}{\mathbb N}\newcommand{\Z}{\mathbb Z}\newcommand{\E}{\mathbb E}
\newcommand{\ip}[2]{\langle #1,#2\rangle}\newcommand{\norm}[1]{\lVert #1\rVert}
\newcommand{\epsc}{\varepsilon_{\mathrm c}}
\DeclareMathOperator{\osc}{osc}\DeclareMathOperator{\spanop}{span}\DeclareMathOperator{\End}{End}\DeclareMathOperator{\diag}{diag}\DeclareMathOperator{\tr}{tr}\DeclareMathOperator{\rad}{rad}\DeclareMathOperator{\TV}{TV}
\title[Input-dependent boundary coding]{Input-Dependent Boundary Coding\\and the Full Error Range}
\author{Qian Qi}
\date{3 October 2026}
\subjclass[2020]{Primary 93B15, 22C05; Secondary 68Q45, 52A27, 60J10}
\begin{document}
\raggedbottom
\begin{abstract}
We determine joint horizon--precision description laws for quantum
instruments that measure a fixed input basis and prepare conditional output
states. Outcome probabilities may depend on the input, and conditional
states may be singular and have changing supports. If the fixed public
rank bounds are $r_{xy}$, put
$V=\sum_x(\sum_y r_{xy}(2n-r_{xy})-1)$. The optimal reusable description
length is $(V/2)\log_2N+V\log_2(1/\delta)+O(1)$, uniformly for $N\ge1$
and $0<\delta\le\delta_*<2$. Error is the unhalved final-state trace
norm over common adaptive quantum testers with references, feedback and
bounded stopping. Rational decoded centres preserve zero conditional
blocks and never increase their ranks. A data-processing retraction
identifies the optimal centre class. A common finite-dimensional
estimator and a covering-multiplicity argument extend the interior and
preparation description laws to the same submaximal error range,
without a two-point separation restriction. We retain the distinct
coherent/preparation law and its narrower stated error range. Dimensions,
rank bounds and the error cap are fixed; no uniformity as the cap tends
to two is asserted. The encoder receives matrix data. Reusable payload,
encoder workspace, learning and physical classical simulation are
different resources.
\end{abstract}
\maketitle
\noindent\textit{Operational interfaces.}
Description compilation returns a mathematical quantum instrument;
numerical streaming processes fully specified matrix data. Neither is a
physical classical device acting on an unknown quantum input, and
neither is merely a transcript-only simulation task. The new code lengths
charge the reusable description, with fixed public dimensions, rank bounds where applicable, horizon and error.
All joint-state errors
are unhalved trace norms; transcript total variation is half the classical
block trace norm.

\section{Introduction}
A quantum instrument can be described once and then used repeatedly.
We ask how many classical description bits suffice to approximate
all interactions with that instrument, uniformly over a common
causal quantum tester. The description specifies a genuine quantum
instrument; it does not replace that instrument by a classical
physical device. The target is supplied to the encoder, rather than
learned from queries.


The principal additional family in the present article is
\[
 \Phi_{\boldsymbol\sigma,y}(X)
   =\sum_x\langle e_x,Xe_x\rangle\sigma_{xy}.
\]
Here a fixed input basis is measured, so outcomes can depend on the
input and the operation destroys input coherence. For public conditional
output-rank bounds $r_{xy}$, let
$V=\sum_x(\sum_y r_{xy}(2n-r_{xy})-1)$. Theorem~\ref{thm:cqentropy71}
proves
\[
 \log_2\mathcal M_N(\delta;\mathcal F_{\boldsymbol r})
 =\frac V2\log_2N+V\log_2(1/\delta)+O(1),
 \qquad 0<\delta\le\delta_*<2.
\]
The remainder is uniform in $N\ge1$ and positive $\delta$ for fixed
dimensions, rank bounds and cap $\delta_*$. No positive eigenvalue
margin is needed. The rational code uses one singular factor code per
input row; a common quantum programme converts its row errors to a
bound for all adaptive testers. The programme is a proof comparison,
not a classical implementation of an unknown quantum input.
Proposition~\ref{prop:inputretraction71} also retracts arbitrary legal
centres to the fixed-readout family without increasing the covering
radius; this retraction does not preserve target rank bounds.

The error cap in this result is obtained by a second argument, not by
extrapolating a local binary test. One informationally complete tester
estimates every target with mean-square error $O(N^{-1})$. Disjoint
estimator events then bound the number of targets assigned to each
covering centre, even when one centre covers more than one point.
Theorem~\ref{thm:fullpacking71} supplies a lower bound at every fixed
unhalved error below two. Combined with the local precision converse,
Corollary~\ref{cor:fullrange71} extends the general positive-margin and
preparation-family laws to $0<\delta\le\delta_*<2$. The endpoint two
is different: one centre covers every instrument there. This estimator
is a common lower-bound witness, not a new optimal-learning theorem.

A complementary coherent boundary family sends
$X$ to $(UXU^*\otimes\sigma_y)_y$. It retains the input quantum data
unitarily, rather than erasing it, and permits all prescribed singular
preparation supports. With $u=d^2-1$ and
$v=\sum_y r_y(2n-r_y)-1$, Theorem~\ref{thm:coherententropy70} proves
\[
 \log_2\mathcal C_N(\delta;\boldsymbol r)
 =(u+v/2)\log_2N+(u+v)\log_2(1/\delta)+O(1).
\]
The finite remainder is uniform in $N\ge1$ and $0<\delta\le1/32$, for
fixed dimensions and public rank bounds. It holds even when covering centres
are arbitrary legal memoryless instruments. Coherent directions have
resolution $\delta/N$, while the preparation directions have resolution
$\delta/\sqrt N$. The coefficient is therefore not half the total
parameter dimension. The upper code specifies a rational unitary and
positive preparation blocks. It is a description of a quantum operation,
not a classical physical implementation of its action.

The proof keeps the two lower witnesses separate and combines their
packings. It gives no simultaneous identification measurement. The upper
bound composes a finite rational unitary atlas with the rank-bounded factor
code. Corollary~\ref{cor:qubitcodec70} makes the qubit construction an
explicit exact-arithmetic encoder. Theorem~\ref{thm:retraction70} settles
a different question: for preparation targets, arbitrary memoryless
covering centres can be retracted exactly to preparation centres. This
retraction need not preserve the target rank bound.

For an $m$-outcome instrument from $M_d$ to $M_n$, write
$s=d^2(mn^2-1)$. Fix $d,n,m$ and a blockwise Choi eigenvalue margin
$0<a<1/(mn)$. Theorem~\ref{thm:jointcode68} proves
\begin{equation*}
 \log_2\mathcal M_N(\delta;a)
 =\frac{s}{2}\log_2 N+s\log_2(1/\delta)+O_{d,n,m,a}(1),
 \qquad N\ge1,\quad 0<\delta\le1/32.
\end{equation*}
The metric is the unhalved final-state trace distance, maximized over
common causal testers with entangled references, arbitrary finite
quantum memory, feedback and bounded public stopping. The same
memoryless instrument is reused at every invocation. The remainder
is independent of both $N$ and $\delta$.

A complementary theorem includes a singular boundary. An
outcome-retaining preparation instrument sends $X$ to
$(\tr(X)\sigma_y)_y$, where $\sigma_y\ge0$ and
$\sum_y\tr\sigma_y=1$. For fixed rank bounds $r_y$, put
\[
                  v=\sum_y r_y(2n-r_y)-1.
\]
Theorem~\ref{thm:rankentropy68} gives
\[
 \log_2\mathcal P_N(\delta;\mathbf r)
 =\frac v2\log_2N+v\log_2(1/\delta)+O_{n,m,\mathbf r}(1)
\]
on the same joint range, including zero outcomes and rank-deficient
blocks, without any eigenvalue floor. The rational code of
Theorem~\ref{thm:factorcode68} preserves zero outcomes and never
increases any block rank. Its integer square roots encode ratios of
triangular-factor coordinates; the encoder need not store algebraic
numbers. The singleton case $v=0$ has zero payload length.

The joint converse rests on a binary product-testing bound that is
linear in the local separation down to arbitrarily small errors.
An affine packing then retains the term $s\log_2(1/\delta)$, which
would be hidden by a fixed-error remainder. On a positive Choi body,
the upper bound uses a two-point classical programme. For preparation
instruments, input erasure instead reduces the full adaptive metric
exactly to the trace distance between product states. Purification
and rank-preserving factor coordinates give a boundary-uniform upper
code. Volume, binary testing, programme arguments and purification
are classical ingredients; their roles are separated in the proofs.

The single-use intrinsic entropy and coordinate codec, and the
fixed-error adaptive law, are retained as complete preliminary
results. No action spectral gap is used in these description theorems.
The complete research edition separately retains the spectral width,
positive numerical streaming and structural realization theory, with
all their original hypotheses and proofs. Neither that larger theory
nor its historical archive is a prerequisite for the present article.

Dimensions, named outcomes, and the requested public parameters are
fixed in the leading payload laws. A rank bound is also public where
used; the finite pivot mask is charged in the explicit factor-code
bound. Leading zero padding is counted conceptually in a fixed-length
code even though the hexadecimal interchange format is canonical and
unpadded. Headers, validation workspace and decoded matrix storage
are additional resources when their sizes vary. Decoding a word can
verify legality, not closeness to an unspecified target. The classical
instrument codec uses its validated encoder image, not its whole
mixed-radix cube, as a legal codebook.

\section{Rational instruments stable under external references}\label{sec:reference66}
A rounding operation on density matrices need not be affine. In
particular, a nonlinear state repair cannot be tensored with the
identity on an unknown reference system. We therefore round a
\emph{description of the instrument}, once and independently of its
input, rather than round the state on which it acts. The resulting
maps are completely positive and jointly trace preserving. All
constants below are explicit in the input and output dimensions and
the number of outcomes.

\subsection{An integral repair of Choi data}
Write $M_d$ for the complex matrices of order $d$. An instrument from
$M_d$ to $M_n$ with $m$ outcomes is a family of completely positive
maps $\Phi_y$ whose sum is trace preserving. We use the input-first,
\emph{unnormalized} Choi convention
\begin{equation}\label{eq:choiconvention66}
 J_y=\sum_{i,j=1}^d |i\rangle\langle j|\otimes
                     \Phi_y(|i\rangle\langle j|),\qquad
 J_y\ge0,\qquad \sum_{y=1}^m\operatorname{Tr}_{n}J_y=I_d.
\end{equation}
Here $\operatorname{Tr}_{n}$ traces the output factor; in particular
$\sum_y\tr J_y=d$, not one. These are the classical Choi criteria
\cite[Theorems~2.22 and~2.26]{Watrous65}. The channel with the outcome
retained is $\mathcal F(X)=\bigoplus_y\Phi_y(X)$. We write $\|\cdot\|_\diamond$
for the completely bounded trace norm.

\begin{lemma}[Choi trace norm and reference stability]\label{lem:choidiamond66}
For any linear map $\Lambda:M_d\to M_n$ with unnormalized Choi matrix
$J(\Lambda)$,
\[
 d^{-1}\|J(\Lambda)\|_1\le\|\Lambda\|_\diamond
                                      \le\|J(\Lambda)\|_1.
\]
Consequently the diamond distance of two instruments, with the outcome
retained, is at most the sum of the trace norms of their Choi-block
differences.
\end{lemma}
\begin{proof}
Applying $\operatorname{id}_d\otimes\Lambda$ to the normalized
maximally entangled state gives the lower bound. For the upper bound,
let $u,v$ be unit vectors in $\C^r\otimes\C^d$. With
$\Omega=\sum_i e_i\otimes e_i$, write
$u=(A\otimes I_d)\Omega$ and $v=(B\otimes I_d)\Omega$.
The matrices $A,B:\C^d\to\C^r$ have Hilbert--Schmidt norm one and
therefore operator norm at most one. Thus
\[
 (\operatorname{id}_r\otimes\Lambda)(|u\rangle\langle v|)
       =(A\otimes I_n)J(\Lambda)(B^*\otimes I_n)
\]
has trace norm at most $\|J(\Lambda)\|_1$. The trace-norm unit ball
is the convex hull of rank-one partial isometries (including zero as a
convex combination of opposite such matrices). A singular-value
decomposition of an arbitrary input matrix therefore proves the same induced
trace-norm bound, independently of $r$. Taking the supremum over $r$
proves the assertion. For an instrument the Choi matrix is, up to a
fixed permutation of factors, the direct sum of its blocks.
\end{proof}

Let $B$ be a positive integer. The supplied approximations $T_y$ are
Hermitian matrices with entries in $B^{-1}\Z[i]$. Assume that every
independent real diagonal coordinate, and every independent real and
imaginary off-diagonal coordinate, differs from that of $J_y$ by at
most $B^{-1}$. The $T_y$ need not be positive or satisfy the marginal
constraint. Put
\begin{align}
 R&=I_d-\sum_y\operatorname{Tr}_{n}T_y,&
 Q_y&=T_y+\mathbf1_{\{y=1\}}R\otimes|e_1\rangle\langle e_1|,
                                                    \label{eq:choirepair66}\\
 c&=2nd(1+mn),&D&=B+mnc,\qquad
 \widetilde J_y=\frac{BQ_y+cI_{dn}}{D}.\label{eq:choigrid66}
\end{align}
The choice of the first outcome and first output coordinate is merely
a fixed convention. Any fixed outcome and output basis coordinate give
the same legality and uniform error bound. If specified zero outcomes
must remain exactly zero, remove them before applying the theorem and
restore them afterwards; $m$ then counts only the repaired outcomes.

\begin{theorem}[Reference-stable rational instrument rounding]\label{thm:choiround66}
For every $d,n,m,B\ge1$, the matrices in \eqref{eq:choigrid66} are
Choi matrices of an instrument. Their numerators are positive semidefinite matrices with Gaussian-integer entries, their common denominator is $D$, and
\begin{equation}\label{eq:choierror66}
 \|\mathcal F-\widetilde{\mathcal F}\|_\diamond
 \le\sum_y\|J_y-\widetilde J_y\|_1
 \le\frac{3mndc}{B+mnc}\le\frac{K(d,n,m)}B,
 \qquad K(d,n,m)=6mn^2d^2(1+mn).
\end{equation}
The construction uses integer additions, multiplication by specified
integers, and the supplied grid entries. It uses no spectral
factorization, positivity projection, or rational Kraus
factorization. Its complete output description has
\[
 O\bigl(m d^2n^2\log(B+mnc+1)\bigr)
\]
bits, with an absolute implicit constant for the fixed encoding.
\end{theorem}
\begin{proof}
Set $E_y=T_y-J_y$. Counting the real coordinates gives
$\|E_y\|_{\rm op}\le\|E_y\|_F\le2nd/B$.
For a Hermitian $E$, the inequalities
$-\|E\|_{\rm op}I\le E\le\|E\|_{\rm op}I$ imply
$\|\operatorname{Tr}_{n}E\|_{\rm op}\le n\|E\|_{\rm op}$.
It follows that
\[
 \|R\|_{\rm op}\le2mn^2d/B,\qquad
             \|Q_y-J_y\|_{\rm op}\le c/B.
\]
Since $J_y\ge0$, the numerator $BQ_y+cI_{dn}$ is positive.
All its entries are Gaussian integers. The partial trace constraint is
exact, because
\[
 \sum_y\operatorname{Tr}_{n}(BQ_y+cI_{dn})
                       =(B+mnc)I_d.
\]
The Choi criteria therefore give a completely positive instrument,
including when some of the original outcomes are zero.

Let $F_y=Q_y-J_y$. The identity
\[
 D(\widetilde J_y-J_y)=BF_y+cI_{dn}-mncJ_y
\]
and $\|F_y\|_1\le dn c/B$ give
\[
 D\sum_y\|\widetilde J_y-J_y\|_1
 \le mndc+mndc+mnc\sum_y\tr J_y=3mndc.
\]
Lemma~\ref{lem:choidiamond66} proves the norm assertion. Finally each
diagonal entry of a positive numerator is at most $D$, by the summed
partial trace constraint. Positivity bounds every off-diagonal modulus
by the geometric mean of two diagonal entries, hence also by $D$.
Encoding $m(dn)^2$ entries and the common denominator proves the
bit bound. All intermediate entries have polynomially bounded bit
length in the stated parameters and the supplied data.
\end{proof}

For exact rational input, directed truncation supplies $T_y$ by integer
division. For uncertain real input, the hypothesis is a certified error
bound on the \emph{final grid coordinates}; an uncertified floating
estimate is not a premise of the theorem. Although $B$ may be a power
of two, $D=B+mnc$ generally is not. The output is rational, not in
general dyadic. The construction may replace a zero outcome by a
small nonzero one; its mass is charged in \eqref{eq:choierror66}.
Known zero outcomes can instead be removed before rounding and restored
as zero blocks afterwards, with $m$ counting the remaining outcomes.

\begin{corollary}[A finite-description instrument net]\label{cor:instrumentnet66}
For fixed $d,n,m$ and $0<\delta\le1$, the preceding construction yields
a finite diamond-$\delta$ net of legal instruments. Each member has
an integer description of length
\[
 O\bigl(m d^2n^2\log(2+K(d,n,m)/\delta+mnc)\bigr).
\]
This is a deterministic description net, not a statistically learned net.
A particular member is computable in polynomial bit time from rational
input data and the grid precision; enumeration of the entire net is
not required.
\end{corollary}
\begin{proof}
Every Choi diagonal is at most one by \eqref{eq:choiconvention66}, and
positivity bounds every entry modulus by one. There are finitely many
possible directed grid truncations at $B=\lceil K/\delta\rceil$.
Apply the theorem to those arising from valid instruments. Computation
of one truncation and its repair uses a polynomial number of elementary
integer operations on polynomial-length operands.
\end{proof}
The CP/TP constraints and adding a scalar identity buffer are not new
principles. Projected least-squares process tomography already repairs
Choi estimates and uses identity mixing to finish a TP-preserving
positivity correction \cite[Section~4]{PLS66}. Here the particular
integer formula supplies an explicit denominator and a deterministic
diamond-error budget from coordinate errors. No nearest-projection,
statistical sample-complexity, or rank-optimality claim is made.

\subsection{Adaptive use and stopped observations}
A tester may retain an arbitrary finite-dimensional quantum memory,
start with a state entangled across that memory and the instrument
input, and insert arbitrary channels and finite-outcome measurements between uses.
The command port is classical: its value may be selected from the
public history and measured tester memory. Coherent access to a
superposition of commands is not an interface assumed here. The
\emph{same} tester is used in the two compared experiments.

\begin{theorem}[Adaptive reference-stable precision budget]\label{thm:adaptivediamond66}
Suppose that at use $t$, uniformly over permitted commands and public
histories,
$\|\mathcal F_{t,a}-\widetilde{\mathcal F}_{t,a}\|_\diamond\le\delta_t$.
For every common tester with a classical command port (not coherent
superpositions of commands), every initial state, and every public stopping
rule $\tau\le N$, the final states, with the stopped history and tester
memory retained, satisfy
\begin{equation}\label{eq:testererror66}
 \|\Omega_\tau-\widetilde\Omega_\tau\|_1
                  \le\sum_{t=1}^N\delta_t.
\end{equation}
The bound is independent of the tester-memory dimension. In particular,
rounding each command by Theorem~\ref{thm:choiround66} with
$B\ge N K(d,n,m)2^L$ gives stopped joint-state error at most $2^{-L}$
and transcript total-variation error at most $2^{-L-1}$.
\end{theorem}
\begin{proof}
Replace the instruments one use at a time. Just before the replaced
use, both members of a neighboring hybrid pair have the same joint
state. Decomposing a classical command/history register into blocks,
the diamond bound applies to each positive block and is multiplied by
its trace. These traces sum to one. The error introduced at this use
is therefore at most $\delta_t$, including any quantum reference.
All subsequent operations in the hybrid pair are the same channels;
complete trace-norm contraction bounds their final difference by the
same quantity. Summing the $N$ hybrid differences proves
\eqref{eq:testererror66}. For a stopping rule retain an absorbing halt
flag, use the identity on halted blocks, and finish all experiments at
use $N$. A public history-dependent stopping rule is then part of the
same tester. The same argument applies, without conditioning on a
possibly rare stopping history. Taking traces and dividing by two
gives the total-variation assertion.
\end{proof}
The hybrid principle is classical; see
\cite[Section~3.3, equation~(3.306)]{Watrous65} and the operational
strategy norms of \cite{Gutoski66}. The conclusion here is a uniform
finite-description approximation by \emph{genuine instruments}. It is
not a classical device that accepts an unknown physical entangled
state and prepares its output. A numerical simulation of a specified
joint state must charge the joint matrix dimension. Nor does the
nonlinear state repair in the preceding section become a channel by
this theorem.

\begin{proposition}[Posterior control at a stopping time]\label{prop:posterior66}
Let $A_h=p_h\rho_h$ and $\widehat A_h=q_h\widehat\rho_h$ be the
positive blocks at a common bounded stopping time, with both total
traces one. At zero simulated mass assign any density matrix to
$\widehat\rho_h$. If
$\sum_h\|A_h-\widehat A_h\|_1\le\delta$, then
\begin{equation}\label{eq:posterior66}
 \sum_h p_h\|\rho_h-\widehat\rho_h\|_1\le2\delta,\qquad
 \Pr_p\{\|\rho_h-\widehat\rho_h\|_1>\eta\}
                              \le\min\{1,2\delta/\eta\}.
\end{equation}
The tail estimate holds for every $\eta>0$. For any event of true probability $\alpha>0$, the two pooled
conditional states have distance at most $\min\{2,2\delta/\alpha\}$,
with an arbitrary conditional state when the simulated event has
probability zero.
\end{proposition}
\begin{proof}
Writing $e_h=\|A_h-\widehat A_h\|_1$, trace contraction gives
$|p_h-q_h|\le e_h$. For positive $p_h$,
\[
 p_h\|\rho_h-\widehat\rho_h\|_1
 \le e_h+|p_h-q_h|\le2e_h.
\]
The inequality remains valid if $q_h=0$. Sum and apply Markov's
inequality. Applying the same calculation to the sums over an event
gives the final assertion; two density matrices have distance at most
two. Zero true mass makes no contribution.
\end{proof}
The proposition gives weighted, tail and pooled-event posterior control,
not simultaneous accuracy of every normalized rare history.

\section{Exact validation of rational descriptions}
The rational codes below require an exact positivity test. We record
the arithmetic fact in the form used by both encoders.

\begin{lemma}[Reduced Schur entries and bit length]\label{lem:schurminors67}
Let $A$ be a Hermitian matrix of order $q$ with Gaussian-integer entries,
whose real and imaginary parts have absolute value at most $2^h$.
In exact Schur elimination, let $I$ be the ordered indices of the
positive pivots already eliminated. Every remaining entry is
\begin{equation}\label{eq:borderminor67}
 S_{ij}=\frac{\det A[I\mathbin{\|}i,I\mathbin{\|}j]}{\det A[I,I]},
\end{equation}
where $I\mathbin{\|}i$ appends $i$ in that order, and the empty
principal determinant is one. Reduced numerators and denominators
have $O(q(h+\log(q+1)))$ bits. Hermitian positive semidefiniteness is
therefore decidable by this elimination in polynomial bit time and
space in the explicit input length. No optimized polynomial exponent
is asserted.
\end{lemma}
\begin{proof}
When the leading pivot $t$ is positive, block congruence reduces
\[
 \begin{pmatrix}t&b^*\\ b&C\end{pmatrix}\ge0
 \quad\Longleftrightarrow\quad C-bb^*/t\ge0.
\]
A negative pivot rejects. At a zero diagonal, each $2\times2$ principal
minor forces the corresponding row and column to vanish; otherwise
reject, and if they vanish delete that row and column. Positive pivots
make $A[I,I]$ positive definite. The bordered determinant identity
\[
 \det\begin{pmatrix}A[I,I]&A[I,j]\\A[i,I]&A_{ij}\end{pmatrix}
 =\det A[I,I]\,(A_{ij}-A[i,I]A[I,I]^{-1}A[I,j])
\]
proves \eqref{eq:borderminor67}, including after zero rows are skipped.
The specified ordering prevents a spurious permutation sign.

A minor of order $k\le q$ has determinant modulus at most
$k^{k/2}(\sqrt2\,2^h)^k$, by the product of its row norms.
Its real and imaginary parts are integers with
$O(k(h+\log(k+1)))$ bits. The principal denominator is a positive
integer. Cancelling common factors cannot enlarge these lengths.
Each arithmetic operation on reduced Gaussian-rational entries thus
has polynomial operand size. There are polynomially many operations
and at most $q^2$ stored entries; ordinary multiplication, division and
the Euclidean algorithm give the stated bit bounds. For rational
rather than integral input, first clear a common positive denominator.
Its bit length is at most the sum of the encoded denominator lengths,
so this preliminary step also has polynomial cost.
\end{proof}

\section{Intrinsic description entropy of instruments}\label{sec:description67}
We now determine the high-resolution description exponent, rather than
only an upper bound on the number of stored matrix entries. The objects
being described remain genuine quantum instruments. A binary description
of such an object is not a physical classical implementation of it.

Fix positive integers $d,n,m$, put $q=dn$, and use the input-first,
unnormalized Choi convention of Section~\ref{sec:reference66}. Let
\[
 \mathcal C=\mathcal C_{d,n,m}
 =\{(J_y)_{y=1}^m:J_y\in\operatorname{Herm}(q),\ J_y\ge0,
                       \ \sum_y\operatorname{Tr}_n J_y=I_d\}.
\]
The corresponding outcome-retaining channel is denoted by $\mathcal F_J$.
On tuples of Hermitian matrices write
\[
 |X|_2=\left(\sum_y\|X_y\|_F^2\right)^{1/2},\qquad
 \|X\|_{\diamond}=\|\mathcal F_X\|_{\diamond}.
\]
Here $\mathcal F_X$ is the linear map defined by the Choi blocks, whether
or not they are positive. Distances in this section are unhalved norms.
The instrument diamond norm means exactly the diamond norm of this
flagged direct-sum channel, not the sum of the individual outcome-map
diamond norms. All $m$ labeled outcome slots are part of the body.
A legal $\eta$-description code is a finite set of codewords with one
member of $\mathcal C$ assigned to each codeword, such that every target
is within diamond distance $\eta$ of a decoded member. Each codeword
selects one fixed memoryless instrument, not a target-dependent random
mixture with uncharged mixing probabilities. A rational code
requires the decoded matrices to have rational real and imaginary parts.
Dimensions, the outcome names and the requested tolerance are fixed
public parameters in the code-length statements. Their encodings are
additional data when they vary.

\begin{lemma}[Affine dimension and an interior ball]\label{lem:affine67}
The affine hull of $\mathcal C$ has real dimension
\begin{equation}\label{eq:dimension67}
 s=mq^2-d^2=d^2(mn^2-1).
\end{equation}
Let $V=\{X:\sum_y\operatorname{Tr}_n X_y=0\}$ and
$J^\circ_y=I_q/(mn)$. For $X\in V$,
\begin{equation}\label{eq:normequi67}
 |X|_2\le\sum_y\|X_y\|_1\le d\|X\|_{\diamond},\qquad
 \|X\|_{\diamond}\le\sqrt{mq}\,|X|_2.
\end{equation}
The $\|\cdot\|_{\diamond}$-ball in $V$ of radius
$r_0=(dmn)^{-1}$, translated by $J^\circ$, is contained in $\mathcal C$.
Moreover $\mathcal C-J^\circ$ lies in the diamond ball of radius two.
\end{lemma}
\begin{proof}
The map from $\operatorname{Herm}(q)^m$ to $\operatorname{Herm}(d)$
given by summing output partial traces is surjective: an arbitrary
$H\in\operatorname{Herm}(d)$ is the image of the tuple with first block
$H\otimes|e_1\rangle\langle e_1|$ and other blocks zero. Its kernel has
codimension $d^2$. All blocks of $J^\circ$ are positive definite, so
small perturbations in that kernel remain legal. This proves
\eqref{eq:dimension67} and that $V$ is the translation space of the
full affine hull, not merely of a containing affine space.

The first two inequalities follow from Lemma~\ref{lem:choidiamond66}
and the direct-sum Choi matrix. The last follows from that lemma,
$\|X_y\|_1\le\sqrt q\|X_y\|_F$, and Cauchy--Schwarz over $y$.
If $\|X\|_{\diamond}\le r_0$, every block has operator norm at most
$|X|_2\le1/(mn)$, which proves positivity of $J^\circ+X$.
The affine constraint is unchanged. Finally a channel has diamond norm
one, so the distance between any two outcome-retaining channels is at
most two. This also follows from complete trace-norm contractivity of
channels, as in \cite[Corollary~3.40]{Watrous65}.
\end{proof}
The lower Choi estimate from Section~\ref{sec:reference66} thus has a
specific converse use: it transfers a diamond ball into the positive
Choi body. Neither this step nor the dimension count uses a spectral gap.

\begin{theorem}[Sharp high-resolution description exponent]\label{thm:entropy67}
Suppose $s>0$. Let $\mathcal N_{\diamond}(\eta)$ and
$\mathcal N^{\mathbb Q}_{\diamond}(\eta)$ be the smallest cardinalities
of legal and rational legal $\eta$-description codes for $\mathcal C$.
For $0<\eta\le1$,
\begin{equation}\label{eq:cover67}
 \max\{1,(r_0/\eta)^s\}
 \le\mathcal N_{\diamond}(\eta)
 \le\mathcal N^{\mathbb Q}_{\diamond}(\eta)
 \le(1+8/\eta)^s.
\end{equation}
In particular, for fixed $d,n,m$, both optimal fixed-length binary
code lengths are
\begin{equation}\label{eq:bits67}
 s\log_2(1/\eta)+O_{d,n,m}(1)\qquad(\eta\downarrow0).
\end{equation}
When $s=0$, $m=n=1$, the unique instrument is the trace map and the
code length is zero.
\end{theorem}
\begin{proof}
Equip $V$ with Lebesgue measure and the diamond norm restricted to $V$.
By Lemma~\ref{lem:affine67}, the target set contains a ball of radius
$r_0$. A union of $M$ balls of radius $\eta$ covering this ball has
volume at most $M\eta^s$ times the unit-ball volume. Hence
$M\ge(r_0/\eta)^s$. Centres anywhere in the affine hull give the same
lower bound.

For the upper bound, take a maximal set of points of $\mathcal C$
whose mutual distances are greater than $\eta/2$. Balls of radius
$\eta/4$ about them are disjoint up to null boundaries and lie in the
ball of radius $2+\eta/4$ about $J^\circ$. The number of points is
therefore at most $(1+8/\eta)^s$, and maximality makes them an
$\eta/2$-net. Rational legal instruments are dense, by the explicit
construction in Theorem~\ref{thm:intrinsiccodec67} below. Replace each
centre by a rational legal point within $\eta/2$. This yields a rational
$\eta$-net of the same cardinality. This dependence is not circular:
the construction below is proved directly from the Choi constraints.

A codebook with $M$ words can be indexed by $\lceil\log_2M\rceil$ bits,
and a code of at most $b$ fixed-length bits has at most $2^b$ words.
Taking logarithms in \eqref{eq:cover67} proves \eqref{eq:bits67}.
The singleton case follows directly from the constraint $J_1=I_d$.
\end{proof}
This is an elementary volumetric theorem for the precise legal Choi
body and norm. The volume comparison is not a new general covering
principle. The coefficient $s$ is sharp in the high-resolution limit
with dimensions fixed; the estimate does not assert optimal constants
in a simultaneous dimension--accuracy limit.

\subsection{An intrinsic rational coordinate code}
Fix an outcome $y_*$ and an output basis index $\alpha_*$. The
independent real coordinates of a Hermitian matrix are its real diagonal
entries and the real and imaginary parts of its strictly upper-triangular
entries. Omit all $d^2$ such coordinates belonging to the submatrix
\[
 ((J_{y_*})_{i\alpha_*,j\alpha_*})_{1\le i,j\le d}.
\]
There are exactly $s$ remaining real coordinates, in a fixed order.
Every coordinate of a legal target has absolute value at most one:
positivity and the marginal identity bound each diagonal by one, and
$|J_{uv}|^2\le J_{uu}J_{vv}$ bounds the off-diagonal entries.

For $B\in\N$, truncate the $s$ retained coordinates toward zero on
the grid $B^{-1}\Z$, producing digits $z_1,\ldots,z_s\in[-B,B]\cap\Z$.
Fill the omitted submatrix by the exact marginal identity:
\begin{equation}\label{eq:intrinsicfill67}
 (Q_{y_*})_{i\alpha_*,j\alpha_*}
 =\delta_{ij}-
   \sum_{(y,\alpha)\ne(y_*,\alpha_*)}
                         (Q_y)_{i\alpha,j\alpha}.
\end{equation}
All summands on the right were retained. Lower-triangular entries are
Hermitian copies, not additional digits.

\begin{theorem}[Exact intrinsic codec]\label{thm:intrinsiccodec67}
Let
\begin{equation}\label{eq:codeconstants67}
 c=2nd(1+mn),\qquad D=B+mnc,\qquad K=3mndc.
\end{equation}
The matrices
\begin{equation}\label{eq:intrinsicrepair67}
 \widehat J_y=(BQ_y+cI_q)/D
\end{equation}
are positive semidefinite with Gaussian-integer numerators, satisfy
$\sum_y\operatorname{Tr}_n\widehat J_y=I_d$, and obey
\begin{equation}\label{eq:intrinsicerror67}
 \|\mathcal F_J-\mathcal F_{\widehat J}\|_{\diamond}
 \le\sum_y\|J_y-\widehat J_y\|_1\le K/D\le K/B.
\end{equation}
They can be encoded by one integer in $[0,(2B+1)^s)$, using at most
$\lceil\log_2((2B+1)^s)\rceil$ bits. Thus $B=\lceil K/\eta\rceil$
achieves the leading coefficient $s$ in \eqref{eq:bits67} by a direct
coordinate code, without enumerating a net. For rational input,
encoding, decoding and validity testing have polynomial bit complexity
in the explicit input and output lengths. This is not a workspace
bound for a trajectory simulator.
\end{theorem}
\begin{proof}
Non-omitted independent-coordinate errors are at most $1/B$. An omitted
coordinate is reconstructed from at most $mn-1$ errors, so its error
is at most $(mn-1)/B$. In particular every independent real error is
at most $mn/B$. Counting the Hermitian copies gives the safe estimate
\[
 \|Q_y-J_y\|_{\mathrm{op}}\le\|Q_y-J_y\|_F
 \le 2qmn/B\le c/B.
\]
The affine constraint is exact by \eqref{eq:intrinsicfill67}. Therefore
$BQ_y+cI_q\ge0$ and its summed output partial trace is
$(B+mnc)I_d$. With $E_y=Q_y-J_y$,
\[
 D(\widehat J_y-J_y)=BE_y+cI_q-mncJ_y.
\]
The sum of the trace norms of each of the three terms is bounded by
$mndc$, using $\sum_y\operatorname{tr}J_y=d$ for the last term.
Lemma~\ref{lem:choidiamond66} proves \eqref{eq:intrinsicerror67}.

Set $M=2B+1$. The integer with base-$M$ digits $z_j+B$ determines
all retained coordinates, and \eqref{eq:intrinsicfill67} then determines
the omitted coordinates. A fixed number of digits is used, including
leading zeros. Only integers whose reconstructed, buffered blocks are
positive are valid codewords. They can be recognized by the exact
Schur test in Lemma~\ref{lem:schurminors67}. Every code produced from a
legal target is valid by the preceding proof. If a total decoder is
required, assign each invalid word the canonical rational instrument
$J^\circ$; unused binary strings can be treated in the same way.
The valid codebook already covers all targets. In the singleton case
the construction decodes the trace map from the empty string.

For rational data the coordinate digits are obtained by exact signed
integer division. Matrix filling, buffering and base conversion are
finite integer operations with polynomial operand lengths. Exact PSD
validation has polynomial cost by Lemma~\ref{lem:schurminors67}.
No eigenvector computation, rational Kraus factorization, random draw,
statistical sample or floating-point decision is involved. For an
arbitrary real target, the coordinate construction proves existence
and density; it is not an algorithm for recovering unknown real input
from finitely many unverified measurements.
\end{proof}

\begin{remark}[Known zero outcomes and public headers]\label{rem:zerocode67}
If a specified set of outcomes is known to be zero, remove it before
coding and restore it afterwards. With $m'\ge1$ active outcomes the
intrinsic dimension is $d^2(m'n^2-1)$, and $m'$ replaces $m$ in the
constants. If the active set is not a fixed public parameter, its mask
must also be encoded. The generic construction does not preserve zeros
without this operation. The choice of $(y_*,\alpha_*)$ affects neither
the dimension nor the bounds. Headers for dimensions, $B$, outcomes
and the chosen coordinate convention are not part of a fixed-parameter
codeword; a self-contained file carries and charges them separately.
\end{remark}

\section{Description entropy under repeated adaptive use}\label{sec:adaptiveentropy67}
The number of calls changes the precision needed in a reusable
instrument description. On a strictly positive Choi body the relevant
scale is $N^{-1/2}$, not the $N^{-1}$ scale supplied by a general
one-slot hybrid. We prove a finite, multiparameter coding statement
against arbitrary common adaptive quantum testers.

For two instruments $J,L\in\mathcal C_{d,n,m}$, define
\begin{equation}\label{eq:adaptivemetric67}
 d_N(J,L)=\sup_{\mathsf T}
       \|\Omega_J^{\mathsf T}-\Omega_L^{\mathsf T}\|_1.
\end{equation}
Here $\mathsf T$ is the same causal tester in both experiments, uses the
same memoryless instrument at most $N$ times, and may have arbitrary
finite-dimensional quantum memory, initially entangled references,
intermediate quantum channels and measurements, and classical
feedback. Fresh ancillary systems of arbitrary finite dimension are
allowed: each tester has finite-dimensional registers at every slot,
but no common dimension bound is imposed in the supremum.
Outcomes and final memory are retained. A public stopping
time bounded by $N$ is allowed. A changed input dimension can be handled
by a tester channel between calls. There is no coherent superposition
of command ports, no tester informed of the hidden choice $J$ versus
$L$, and no stopping based on inaccessible simulator state. Taking
suprema of trace distances gives the triangle inequality. In
particular $d_N\le2$. Public stopping is included by an absorbing flag,
as in Theorem~\ref{thm:adaptivediamond66}.

For a fixed $a$ with $0<a<1/(mn)$ let
\[
 \mathcal C_a=\{J\in\mathcal C:J_y\ge aI_{dn}\text{ for every }y\}.
\]
This is a Choi eigenvalue margin, not a spectral gap for an action or a
mixing hypothesis on trajectories. Its affine dimension is the same
$s=d^2(mn^2-1)$ as before.

\begin{lemma}[Two-point classical programme bound]\label{lem:programme67}
Suppose $J_y,L_y\ge\lambda I_{dn}$ for every outcome and
$r=|J-L|_2\le\lambda$, where $\lambda>0$. Then, for every $N\ge1$,
\begin{equation}\label{eq:rootN67}
 d_N(J,L)\le\min\{2,\,2\sqrt N\,r/\lambda\}.
\end{equation}
The same bound covers arbitrary initially entangled references and
bounded public stopping in \eqref{eq:adaptivemetric67}.
\end{lemma}
\begin{proof}
There is nothing to prove for $r=0$. Put $K=(J+L)/2$ and
$U=(J-L)/r$. Since $U\in V$ and each $\|U_y\|_{\mathrm{op}}\le1$,
\[
                      K^+=K+\lambda U,\qquad K^-=K-\lambda U
\]
are legal instruments. With
\[
 p_\pm=\tfrac12\pm\frac r{4\lambda},\qquad
 q_\pm=\tfrac12\mp\frac r{4\lambda},
\]
we have $J=p_+K^++p_-K^-$ and $L=q_+K^++q_-K^-$.
All four probabilities are at least $1/4$. The elementary inequality
$\log x\le x-1$ implies
\[
 D(p\|q)\le\sum_{b\in\{+,-\}}\frac{(p_b-q_b)^2}{q_b}
                                      \le\frac{2r^2}{\lambda^2}.
\]
Choose the programme signs independently in advance for all $N$
slots. Conditional on a sign string, a fixed tester acts on exactly
the same sequence of genuine quantum instruments in the two
experiments. It therefore defines one channel from the classical sign
string to the final quantum--classical state. This remains true with
adaptive feedback: conditioning fixes the instruments, not the
tester's outcomes. If the tester stops early, the unused signs are
ignored. Contractivity and the classical product identity give
\[
 \|\Omega_J^{\mathsf T}-\Omega_L^{\mathsf T}\|_1
 \le2\operatorname{TV}(p^{\otimes N},q^{\otimes N})
 \le\sqrt{2N D(p\|q)}\le2\sqrt N\,r/\lambda.
\]
The middle bound is the usual Pinsker inequality. To recall an elementary
proof, coarse-grain two distributions to the set where the first is
larger. The log-sum inequality reduces relative entropy to binary
relative entropy. Its second derivative in the first argument is
$1/[x(1-x)]\ge4$, so integration about its minimum proves
$D\ge2\operatorname{TV}^2$. Taking the supremum over testers proves
the assertion.
\end{proof}
The classical-programme idea is an established method in quantum
metrology; Demkowicz-Dobrza\'nski, Ko\l ody\'nski and Gu\c{t}\u{a}
\cite[Equations~(4)--(6), Methods]{DKG67} use parameter-dependent
mixtures of fixed channels to bound distinguishability. The lemma
above gives the finite two-point trace-norm estimate needed here,
including common adaptive testers and public stopping. It is not a
claim that the mixture method or the square-root scale is a new
metrological principle. The operations $K^\pm$ remain quantum
operations, not classical physical substitutes for unknown input.

\begin{lemma}[Repeated Choi tests]\label{lem:choitesting67}
For any two instruments $J,L\in\mathcal C$, with $t=|J-L|_2$,
\begin{equation}\label{eq:choitesting67}
 d_N(J,L)\ge2-4\exp\left(-\frac{Nt^2}{8d^2}\right).
\end{equation}
A negative right-hand side is simply a vacuous lower bound.
\end{lemma}
\begin{proof}
At each call supply half of a normalized maximally entangled state and
retain its other half and the outcome. The resulting density matrices
are $A=d^{-1}\bigoplus_yJ_y$ and $B=d^{-1}\bigoplus_yL_y$.
They have trace one. The projector on the positive part of $A-B$
gives a binary measurement whose probability gap is
\[
 g=\tfrac12\|A-B\|_1\ge t/(2d).
\]
The same binary measurement is made on each of the $N$ independent
outputs. With a threshold halfway between the two means, each wrong-side
probability is at most $\exp(-Ng^2/2)$. Indeed for a Bernoulli variable
$Z$, the second derivative of
$\log\mathbb E e^{u(Z-\mathbb EZ)}$ is the variance of a tilted
Bernoulli variable and is at most $1/4$. The value and first derivative
at zero vanish; integration and Chernoff's inequality give
$\Pr\{N^{-1}\sum Z_i-\mathbb EZ\ge v\}\le e^{-2Nv^2}$,
and likewise for the other tail. Set $v=g/2$.
The total variation distance between the two product experiments is
at least $1-2e^{-Ng^2/2}$; the unhalved trace norm of their classical
outputs is twice this. This particular nonadaptive tester is allowed
in \eqref{eq:adaptivemetric67}, and the result follows. The tester may
depend on the pair $J,L$, as is permitted when lower-bounding a metric;
it is the same tester for the two members of that pair.
\end{proof}

\begin{theorem}[Sharp horizon exponent for interior descriptions]\label{thm:adaptiveentropy67}
Fix $d,n,m$ with $s>0$, a margin $0<a<1/(mn)$ and a tolerance
$0<\delta<1$. Let $\mathcal M_N(\delta;a)$ be the minimum number of
legal decoded instruments whose $d_N$-balls of radius $\delta$ cover
$\mathcal C_a$. A decoded instrument need not itself have margin $a$.
Write $\mathcal M_N^{\mathbb Q}$ when rational decoded instruments are
required. Set
\[
 R=\frac1{mn}-a,\qquad
 c_\delta=\sqrt{8d^2\log\frac4{1-\delta}},\qquad
 B_N=\left\lceil\frac Ka\left(2+
                         \frac{4\lceil\sqrt N\rceil}{\delta}\right)\right\rceil,
\]
where $K$ is from \eqref{eq:codeconstants67}. Then
\begin{equation}\label{eq:adaptivecover67}
 \max\{1,(R\sqrt N/c_\delta)^s\}
 \le\mathcal M_N(\delta;a)
 \le\mathcal M_N^{\mathbb Q}(\delta;a)
 \le(2B_N+1)^s.
\end{equation}
Consequently the minimum fixed-length description satisfies
\begin{equation}\label{eq:adaptivebits67}
 \log_2\mathcal M_N^{\mathbb Q}(\delta;a)
       =\frac s2\log_2N+O_{d,n,m,a,\delta}(1),
\end{equation}
and the same formula holds without the rationality restriction.
The upper bound is achieved by the exact intrinsic codec applied once
at grid $B_N$; its decoded instrument is reused at every call.
\end{theorem}
\begin{proof}
The tuple-Frobenius ball in $V$ of radius $R$ about $J^\circ$ is
contained in $\mathcal C_a$. Choose a maximal set of points in that
ball with pairwise distances greater than $c_\delta/\sqrt N$.
Balls of that radius about the selected points cover the original
ball. Euclidean volume comparison therefore gives at least
$(R\sqrt N/c_\delta)^s$ points. By Lemma~\ref{lem:choitesting67}, their
pairwise $d_N$ distances are greater than
\[
 2-4e^{-c_\delta^2/(8d^2)}=1+\delta>2\delta.
\]
One $d_N$-ball of radius $\delta$ cannot cover two of these points,
regardless of the location of its legal centre. This proves the lower
bound.

For the upper bound use Theorem~\ref{thm:intrinsiccodec67}. It gives
\[
 r=|J-\widehat J|_2\le\sum_y\|J_y-\widehat J_y\|_1
 \le K/B_N\le\min\{a/2,\ a\delta/(4\sqrt N)\}.
\]
Every block of $\widehat J$ is consequently at least $(a/2)I$, and
both $J$ and $\widehat J$ satisfy Lemma~\ref{lem:programme67} with
$\lambda=a/2$ and $r\le\lambda$. Its conclusion gives
$d_N(J,\widehat J)\le4\sqrt N\,r/a\le\delta$.
There are at most $(2B_N+1)^s$ possible digit strings. All codewords
arising from targets in $\mathcal C_a$ decode legally. The remaining
codewords may be rejected or assigned $J^\circ$; this does not affect
the cover. Since $B_N=O_{d,n,m,a,\delta}(\sqrt N)$, taking logarithms
in \eqref{eq:adaptivecover67} proves \eqref{eq:adaptivebits67}.
\end{proof}

\begin{remark}[What is charged]\label{rem:adaptivescope67}
The lower bound concerns the number of possible reusable mathematical
instrument descriptions, not the mutable memory of a simulator that
receives and stores the entire instrument as input. It applies to all
full-dimensional strictly positive instrument families above, without
an expanding unitary subsystem. It therefore complements, rather than
replaces, the hypothesis-sensitive streaming-space lower bounds.
For variable $a,\delta,d,n,m$, the displayed finite inequalities remain
valid; \eqref{eq:adaptivebits67} holds with those parameters fixed.
The estimates of this section alone assert neither an optimal uniform
dimension constant nor a complete joint $N$--$\delta$ asymptotic.
The latter is supplied, for fixed dimensions and margin, by
Theorem~\ref{thm:jointcode68}. The rational implementation takes
rational $a,\delta$ and an asserted margin, which it can verify exactly.
No samples of an unknown channel are supplied or required.
\end{remark}

\begin{proposition}[The boundary is quantitatively different]\label{prop:boundaryphase67}
There is no constant $C$ such that
$d_N(J,L)\le C\sqrt N\,|J-L|_2$ holds for every pair of qubit channels
and every $N$. In particular the strict margin in
Theorem~\ref{thm:adaptiveentropy67} cannot be removed from its proof by
setting $a=0$.
\end{proposition}
\begin{proof}
Let $U_\theta=\operatorname{diag}(1,e^{i\theta})$ and compare its
unitary channel to the identity channel. For their unnormalized Choi
matrices,
\[
 |J(U_\theta)-J(I)|_2
       =2\sqrt2\,|\sin(\theta/2)|\le\sqrt2\,|\theta|.
\]
At $\theta=\pi/N$, sequentially applying the unknown channel $N$
times to $|+\rangle$ gives respectively $|-\rangle$ and $|+\rangle$.
The final trace distance is two. The proposed upper bound would tend
to zero. The general additive diamond hybrid remains valid at this
boundary; the proposition asserts no global covering exponent for
boundary or rank-constrained instrument classes.
\end{proof}

\subsection{Description, learning and physical simulation}
A reusable rational description is different from a protocol that
replaces quantum communication by classical messages.
Naik, Zartab, Gisin and Banik \cite[Theorems~2 and~5]{NZGB67} show a
finite-classical-message obstruction for the perfect qubit channel in
a reference-sensitive joint-measurement model, and finite-resource
possibility for noisy depolarizing channels. Their receiver retains a
quantum system; the task is physical simulation. Here the decoded
objects and the vertices in Lemma~\ref{lem:programme67} are genuine
quantum instruments. Our numerical algorithms instead start from
specified matrices. Neither construction supplies a classical physical
device accepting unknown entangled input.

Unknown-channel learning charges channel queries used to infer a
classical description. Mele and Bittel \cite[Theorem~I.1]{MB67}
prove dimension-optimal query complexity at fixed diamond accuracy;
Oufkir and Girardi \cite{OG67} give further accuracy-dependent query
lower bounds. The input here is already specified, and the lower bound
in \eqref{eq:adaptivecover67} counts descriptions, not the queries
needed to select one from unknown data. The elementary packing and
testing ingredients are shared methods, not a priority claim that
those learning theorems are consequences of this code.

\section{Joint horizon and accuracy in description length}\label{sec:jointdescription68}
The exponential test in Lemma~\ref{lem:choitesting67} resolves pairs at
fixed error, but its lower bound is vacuous when the two experiments are
very close. A bounded statistic gives the local estimate needed to retain
the accuracy parameter. Throughout this section $N\ge1$ and the adaptive
distance $d_N$ is the unhalved final-state trace distance maximized over
one common tester, as in \eqref{eq:adaptivemetric67}. The unknown device
is the same memoryless instrument at every call; commands are classical,
and stopping is public and bounded by $N$.

\begin{lemma}[A local binary test]\label{lem:binarylocal68}
For $0\le p,q\le1$,
\begin{equation}\label{eq:binarylocal68}
 \TV\bigl(\operatorname{Bern}(p)^{\otimes N},
          \operatorname{Bern}(q)^{\otimes N}\bigr)
 \ge \frac3{32}\min\{\sqrt N\,|p-q|,1\}.
\end{equation}
The estimate is uniform up to the endpoints of the probability interval.
\end{lemma}
\begin{proof}
Assume $p<q$ and first suppose $q-p\le N^{-1/2}$. Put
$b=N(p+q)/2$ and let $f:\R\to[0,1]$ be the nondecreasing function
which is affine with slope $(8\sqrt N)^{-1}$ on
$[b-4\sqrt N,b+4\sqrt N]$, zero to its left, and one to its right.
For $S\sim\operatorname{Bin}(N,t)$, differentiating the finite binomial
polynomial gives the exact identity
\[
 \frac d{dt}\E_t f(S)
 =N\E_t\{f(T+1)-f(T)\},\qquad T\sim\operatorname{Bin}(N-1,t).
\]
It holds on $(0,1)$ and extends continuously to both endpoints. For
$t\in[p,q]$, the distance from $\E_tT$ to $b$ is at most
$\sqrt N/2+1$. Also $\operatorname{Var}_tT\le(N-1)/4$. Chebyshev's
inequality gives
\[
 \Pr_t\{|T-\E_tT|\le\sqrt N\}\ge3/4.
\]
On this event, both $T$ and $T+1$ belong to the affine interval:
their distance to $b$ is at most $3\sqrt N/2+2\le4\sqrt N$.
Since $f$ is nondecreasing everywhere,
\[
 \frac d{dt}\E_t f(S)\ge\frac3{32}\sqrt N.
\]
Integration from $p$ to $q$, and the characterization of total
variation by $[0,1]$-valued statistics, prove the assertion in this
case. The statistic depends only on the count, which is a function
of the full binary string.

If $q-p>N^{-1/2}$, apply the same construction to
$p$ and $q_0=p+N^{-1/2}<q$. Couple Bernoulli samples by common
uniform random variables. Monotonicity of $f$ then gives
$\E_qf(S)\ge\E_{q_0}f(S)$, and the preceding lower bound is $3/32$.
The cases $p=q$ and $q<p$ follow by equality and symmetry.
\end{proof}
The proof uses only binomial differentiation and a bounded statistic;
it requires neither a normal approximation nor a lower variance bound.
In particular its constant does not deteriorate for a rare outcome.

\begin{corollary}[Local repeated Choi separation]\label{cor:localchoi68}
For the instrument body $\mathcal C_{d,n,m}$, let $t=|J-L|_2$.
Then
\begin{equation}\label{eq:localchoi68}
 d_N(J,L)\ge\frac3{16}
                  \min\left\{\frac{\sqrt N}{2d}t,1\right\}.
\end{equation}
\end{corollary}
\begin{proof}
Use the same normalized Choi inputs as in
Lemma~\ref{lem:choitesting67}. The positive spectral projection of
$d^{-1}\bigoplus_y(J_y-L_y)$ defines a binary test with probability
gap at least $t/(2d)$. Repeating this test independently and applying
Lemma~\ref{lem:binarylocal68} gives \eqref{eq:localchoi68}; classical
trace distance is twice total variation. This is one admissible
nonadaptive tester. For each pair it is the same tester in both
experiments, and it is allowed in the supremum defining $d_N$.
\end{proof}

\begin{theorem}[The joint interior coding law]\label{thm:jointcode68}
Fix $d,n,m$ with $s=d^2(mn^2-1)>0$ and $0<a<1/(mn)$.
For all $N\ge1$ and $0<\delta\le1/32$,
\begin{equation}\label{eq:jointcode68}
 \max\left\{1,\left(\frac{R\sqrt N}{32d\delta}\right)^s\right\}
 \le\mathcal M_N(\delta;a)
 \le\mathcal M_N^{\Q}(\delta;a)\le(2B_N+1)^s,
\end{equation}
where $R=1/(mn)-a$ and $B_N$ is the explicit grid in
Theorem~\ref{thm:adaptiveentropy67}. Consequently the optimal legal
fixed-length description, with or without rational decoded matrices,
has length
\begin{equation}\label{eq:jointbits68}
 \frac s2\log_2N+s\log_2(1/\delta)+O_{d,n,m,a}(1),
\end{equation}
uniformly on this whole range of $N$ and $\delta$. The decoded
instrument is reused unchanged, and no spectral-gap hypothesis occurs.
\end{theorem}
\begin{proof}
The tuple-Frobenius ball of radius $R$ about $J^\circ$ in $V$ is
contained in $\mathcal C_a$. Take a maximal set in this ball whose
pairwise Euclidean distances are greater than
$h=32d\delta/\sqrt N$. Maximality and $s$-dimensional volume give at
least $(R/h)^s$ points. Corollary~\ref{cor:localchoi68} separates
each pair by more than $2\delta$ in $d_N$: before the minimum saturates
its value is greater than $3\delta$, and after saturation it is
$3/16>2\delta$. Thus no legal centre of a $d_N$-ball of radius
$\delta$ can cover two packing points, even if that centre has no
positive eigenvalue margin. This proves the lower bound.

The intrinsic rational codec and its classical-programme estimate in
Theorem~\ref{thm:adaptiveentropy67} give the upper bound with the
same $B_N$. Since $N\ge1$ and $\delta\le1/32$, there is a constant
$C=C(d,n,m,a)$ such that $2B_N+1\le C\sqrt N/\delta$. The lower
bound and this inequality give \eqref{eq:jointbits68}, including
rounding a logarithm up to a fixed-length binary code. When the
volumetric lower expression is less than one, its logarithm is negative;
the nonnegative code length still gives the same uniform additive
estimate. Neither bound charges the queries of a learning procedure.
\end{proof}
The earlier fixed-error theorem is retained. The new assertion is the
uniform additive remainder in \eqref{eq:jointbits68}, which is
independent of both the horizon and the accuracy. It does not assert
sharp dimension-dependent constants, nor is it a lower bound for the
mutable workspace of the variable-description trajectory algorithm.

\section{Boundary-uniform coding of preparation instruments}\label{sec:boundarycode68}
Strict positivity is essential to the two-point programme argument, as
Proposition~\ref{prop:boundaryphase67} shows. Nevertheless it is not
necessary for a complete coding law in an irreversible family. We
consider all outcome-retaining instruments that erase their input and
prepare a specified quantum--classical state. This family contains zero
outcomes and arbitrary rank-deficient output blocks.

For $m,n\ge1$ put
\[
 \mathcal S_{m,n}=\{(\sigma_y)_{y=1}^m:\sigma_y\in\operatorname{Herm}(n),
                  \ \sigma_y\ge0,\ \sum_y\tr\sigma_y=1\}.
\]
Write $\boldsymbol\sigma=\bigoplus_y\sigma_y$. For any input
dimension $d\ge1$, the associated instrument is
\begin{equation}\label{eq:replacer68}
 \mathcal R_{\sigma,y}(X)=\tr(X)\sigma_y,\qquad
 J_y(\mathcal R_\sigma)=I_d\otimes\sigma_y.
\end{equation}
The output label $y$ is retained. Its conditional quantum state is
$\sigma_y/\tr\sigma_y$ when that trace is positive. This is a
specified mathematical quantum instrument; a description of it is not
a numerical algorithm acting on an unknown physical input.

\begin{lemma}[Reduction to a product-state experiment]\label{lem:productreplacer68}
For arbitrary $\sigma,\tau\in\mathcal S_{m,n}$,
\begin{equation}\label{eq:productreplacer68}
 d_N(\mathcal R_\sigma,\mathcal R_\tau)
   =\|\boldsymbol\sigma^{\otimes N}
                -\boldsymbol\tau^{\otimes N}\|_1.
\end{equation}
Here the tester may have arbitrary quantum memory, initial reference
entanglement, feedback and a public stopping time at most $N$.
\end{lemma}
\begin{proof}
Supply $N$ independent program states $\boldsymbol\sigma$ to one
fixed processor. At the next device call the processor discards the
incoming device system and returns the next program state, including
its classical outcome. Between calls it performs the prescribed tester
operations. This reproduces the experiment exactly, even when the
incoming system is entangled with the tester memory: the incoming
system is traced out in \eqref{eq:replacer68}. Unused program states
are ignored after a public stop. The processor is the same channel
for $\sigma$ and $\tau$, so trace-norm contraction gives the upper
bound. Conversely, a tester can make exactly $N$ calls and retain
every outcome and quantum output. Its final state is the displayed
product, which proves equality. The processor is a quantum channel,
not a claimed finite-classical-message physical simulator.
\end{proof}

\subsection{Triangular factors and a rational code}
Every $\sigma\in\mathcal S_{m,n}$ has triangular factors
$L_y$ with $\sigma_y=L_yL_y^*$ and nonnegative real diagonal entries.
The usual no-pivot Cholesky procedure also covers singular matrices:
a zero diagonal pivot forces the entire residual row and column to
vanish and contributes a zero factor column. A block of rank $r$
therefore has exactly $r$ positive pivots. All factors together satisfy
$\sum_y\|L_y\|_F^2=1$.

Fix public bounds $0\le r_y\le n$, not all zero, and write
\begin{equation}\label{eq:rankdimension68}
 D=\sum_{y=1}^m r_y(2n-r_y),\qquad v=D-1,
 \qquad
 \mathcal S_{\boldsymbol r}
   =\{\sigma\in\mathcal S_{m,n}:\operatorname{rank}\sigma_y\le r_y\}.
\end{equation}
If the positive pivot positions of $L_y$ are $T_y\subseteq\{1,\ldots,n\}$,
record only the real diagonal and the real and imaginary coordinates
strictly below each such pivot. Their total number is
\[
 h=\sum_y\sum_{j\in T_y}\{1+2(n-j)\}\le D.
\]
The resulting real coordinate vector $x\in\R^h$ has norm one.
The inequality follows by placing each block's pivots as early as
possible. It holds also when the actual ranks are smaller than the
public bounds.

Choose a coordinate $j$ of maximal absolute value, put $a=|x_j|$,
and set $w=x/a$. For an integer $B\ge1$, take
\begin{equation}\label{eq:factorgrid68}
 z_j=B\operatorname{sgn}(x_j),\qquad
 z_i=\operatorname{sgn}(x_i)\lfloor B|w_i|\rfloor\quad(i\ne j).
\end{equation}
Insert these integers in the recorded factor positions to obtain
Gaussian-integer triangular matrices $Z_y$; all unrecorded columns
are zero. Define
\begin{equation}\label{eq:factoroutput68}
 T=\sum_y\|Z_y\|_F^2,\qquad
 \widehat\sigma_y=Z_yZ_y^*/T.
\end{equation}
Since $|z_j|=B$, the denominator satisfies $B^2\le T\le hB^2$.

\begin{theorem}[A boundary-uniform rational factor code]\label{thm:factorcode68}
The output \eqref{eq:factoroutput68} lies in
$\mathcal S_{\boldsymbol r}$, preserves every zero target outcome,
and does not increase the rank of any individual block. For all
$N\ge1$,
\begin{equation}\label{eq:factorerror68}
 d_N(\mathcal R_\sigma,\mathcal R_{\widehat\sigma})
 \le\min\left\{2,\frac{4\sqrt{N(h-1)}}B\right\}
 \le\min\left\{2,\frac{4\sqrt{Nv}}B\right\}.
\end{equation}
At fixed $m,n,\boldsymbol r,B$ a fixed-length code needs at most
\begin{equation}\label{eq:factorbits68}
 mn+\left\lceil\log_2\bigl(2D(2B+1)^v\bigr)\right\rceil
\end{equation}
bits. For rational target blocks, its encoder and decoder use exact
rational and integer arithmetic, including integer square roots; no
algebraic-number oracle, positive eigenvalue margin, or Kraus
factorization is required.
\end{theorem}
\begin{proof}
The product $Z_yZ_y^*$ is positive semidefinite, and normalization
makes the sum of its traces exactly one. Its number of possibly nonzero
columns is no larger than that of $L_y$, so its rank does not exceed
the target rank. When an outcome is zero it has no recorded columns
and remains exactly zero.

Put $z'=z/B$. Its $j$th coordinate agrees with that of $w$; every
other coordinate differs by less than $1/B$. Both $\|w\|$ and
$\|z'\|$ are at least one. The elementary normalization estimate gives
\begin{equation}\label{eq:factornormalize68}
 \left\|x-\frac z{\|z\|}\right\|
 \le\frac{2\|w-z'\|}{\|w\|}
 \le\frac{2\sqrt{h-1}}B.
\end{equation}
Use the normalized vectors
\[
 |L\rangle=\sum_{y,i,j}(L_y)_{ij}|y,i\rangle|y,j\rangle,
 \qquad
 |\widehat L\rangle=T^{-1/2}\sum_{y,i,j}(Z_y)_{ij}|y,i\rangle|y,j\rangle.
\]
Their reduced states on the first factor are
$\boldsymbol\sigma$ and $\widehat{\boldsymbol\sigma}$. Their vector
distance is the left side of \eqref{eq:factornormalize68}. For any
unit vectors $u,v$, diagonalization on their two-dimensional span gives
\[
 \||u\rangle\langle u|^{\otimes N}
       -|v\rangle\langle v|^{\otimes N}\|_1
   =2\sqrt{1-|\langle u,v\rangle|^{2N}}
   \le2\sqrt N\,\|u-v\|.
\]
Indeed $1-t^N\le N(1-t)$ on $[0,1]$ and
$1-|\langle u,v\rangle|^2\le2(1-\Re\langle u,v\rangle)$.
Taking partial traces and using Lemma~\ref{lem:productreplacer68}
proves \eqref{eq:factorerror68}. These are standard purification and
product-fidelity estimates; see \cite[Section~3.2]{Watrous65}. Their
short proof specifies the normalization used here.

The pivot sets use $mn$ bits. Given them, there are at most $2h$
anchor/sign choices and $(2B+1)^{h-1}$ remaining digit strings.
Pad their index to the maximum length in \eqref{eq:factorbits68}.
An invalid header may be rejected, or sent by a total mathematical
decoder to a fixed legal state. Valid factor words always decode
positively, even if they did not arise from this encoder.

For rational input, maintain the exact rational Schur residual $A$.
At a positive pivot $u=A_{jj}$, the factor diagonal is $\sqrt u$
and a lower entry is $A_{ij}/\sqrt u$. Thus the square of each of
its real coordinates is rational, respectively $u$, or
$(\Re A_{ij})^2/u$, or $(\Im A_{ij})^2/u$; its sign is known.
Updating the residual uses rational arithmetic only. Choose the
largest squared coordinate $b_j>0$ by exact comparisons. If another
coordinate has squared value $b_i$, then
\[
 \lfloor B|w_i|\rfloor
   =\operatorname{isqrt}\!\left(\left\lfloor B^2b_i/b_j\right\rfloor\right).
\]
Consequently no square root of a rational number must be represented.
The ratios-of-minors bounds in Lemma~\ref{lem:schurminors67} give
polynomial operand lengths during this preprocessing. Integer square
root, products and mixed-radix coding have polynomial bit complexity
in the explicit input and output lengths. Encoding workspace and the
expanded decoded matrix are not identified with the compressed payload.
\end{proof}

\subsection{Sharp description laws on the boundary}
Let $\mathcal P_N(\delta;\boldsymbol r)$ be the minimum number of
legal preparation instruments whose $d_N$-balls of radius $\delta$
cover $\mathcal S_{\boldsymbol r}$; the decoded centre may have any
output rank. Add a superscript $\Q$ when rational centres are required.
For an upper construction our centres will obey the same rank bounds.

\begin{theorem}[Rank-stratified joint coding law]\label{thm:rankentropy68}
Fix $m,n,\boldsymbol r$ as in \eqref{eq:rankdimension68}. If $v>0$,
there exist constants $c,C>0$, depending only on these fixed data,
such that, for every $N\ge1$ and $0<\delta\le1/32$,
\begin{equation}\label{eq:rankcover68}
 \max\{1,(c\sqrt N/\delta)^v\}
 \le\mathcal P_N(\delta;\boldsymbol r)
 \le\mathcal P_N^{\Q}(\delta;\boldsymbol r)
 \le C(\sqrt N/\delta)^v.
\end{equation}
The optimal legal and rational fixed-length descriptions therefore have
length
\begin{equation}\label{eq:rankbits68}
 \frac v2\log_2N+v\log_2(1/\delta)+O_{m,n,\boldsymbol r}(1)
\end{equation}
uniformly in $N$ and $\delta$. The statement has no spectral-gap or
positive-eigenvalue-margin hypothesis. It includes every rank-deficient
boundary allowed by the public bounds. If $v=0$, the family is a
singleton and its optimal description length is zero.
\end{theorem}
\begin{proof}
For the upper bound take $B=\lceil4\sqrt{Nv}/\delta\rceil$ in
Theorem~\ref{thm:factorcode68}; for rational implementation it suffices
to replace $\sqrt{Nv}$ by its integer ceiling. Its error is at most
$\delta$, and its codebook cardinality is at most
$2^{mn}2D(2B+1)^v$. This has the asserted upper order uniformly on
the stated parameter range.

For completeness we construct a full-dimensional local patch for the
lower bound. For every block with $r_y>0$, take an $n$ by $r_y$
lower-trapezoidal factor: entries above the first $r_y$ diagonal
positions vanish, and these diagonal entries are positive real
numbers. These factors have $D$ independent real coordinates. Impose
$\sum_y\|L_y\|_F^2=1$ and work near the point with all these diagonal
entries equal to $(\sum_y r_y)^{-1/2}$ and all other entries zero.
An open patch of this unit sphere is a smooth $v$-dimensional manifold.
The map
\[
                         (L_y)_y\longmapsto(L_yL_y^*)_y
\]
is one-to-one there and has a smooth inverse on its image. To see the
inverse explicitly, the leading $r_y$ by $r_y$ principal block is
positive definite and determines its unique positive-diagonal Cholesky
factor $A_y$. The bottom-left block is $B_yA_y^*$, so it determines
$B_y$ by multiplication by $(A_y^*)^{-1}$. These operations are smooth
near the chosen point. When $r_y=0$, the entire block is identically
zero and contributes no coordinate.

Choose a closed coordinate ball of sufficiently small positive radius
inside this patch. The chart and its inverse have bounded derivatives
on a slightly larger neighborhood; hence distances between parameters
in the ball are bounded above and below by constant multiples of the
tuple-Frobenius distance between their states. More explicitly one
may bound the derivative of the inverse on a convex neighborhood of
the chosen positive principal blocks, where the same Cholesky formulas
extend smoothly. A maximal Euclidean packing of this parameter ball
at scale $A\delta/\sqrt N$ therefore has at least
$(c\sqrt N/\delta)^v$ points whose state distances exceed
$32\delta/\sqrt N$, after increasing the fixed $A$.

The binary argument of Corollary~\ref{cor:localchoi68}, applied
directly to the trace-one block states (input dimension one), separates
each pair by more than $2\delta$ after $N$ copies. By
Lemma~\ref{lem:productreplacer68} the same separation holds in $d_N$.
One radius-$\delta$ ball, regardless of the rank of its legal centre,
cannot contain two packing points. This proves the lower bound and
then \eqref{eq:rankbits68}. Finally $D=1$ requires $n=1$ and exactly
one bound $r_y=1$; all other bounds vanish. The sole permitted state
has unit mass on that outcome, as asserted.
\end{proof}

\begin{corollary}[All ranks, with an explicit lower constant]\label{cor:fullboundary68}
For the full family $\mathcal S_{m,n}$ put $v=mn^2-1>0$. Then
\begin{equation}\label{eq:fullboundary68}
 \max\left\{1,\left(\frac{\sqrt N}{64mn\delta}\right)^v\right\}
 \le\mathcal P_N(\delta)\le\mathcal P_N^{\Q}(\delta)
 \le 2^{mn}2mn^2(2B+1)^v,
 \qquad B=\left\lceil\frac{4\sqrt{Nv}}\delta\right\rceil.
\end{equation}
In particular a general $n$-dimensional prepared state has coefficient
$n^2-1$, while the class of pure prepared states has coefficient
$2n-2$ in \eqref{eq:rankbits68}. A classical $m$-outcome law has
coefficient $m-1$, including probabilities equal to zero.
\end{corollary}
\begin{proof}
The upper bound is Theorem~\ref{thm:factorcode68} with every $r_y=n$.
For the lower bound, the affine space of block states has dimension
$mn^2-1$. Its tuple-Frobenius ball of radius $1/(2mn)$ about the
blocks $I_n/(mn)$ lies in $\mathcal S_{m,n}$. Repeat the packing proof
of Theorem~\ref{thm:jointcode68} with input dimension one. For pure
states take $m=1,r_1=1$ in Theorem~\ref{thm:rankentropy68}; for
classical laws take $n=1$.
\end{proof}

\begin{remark}[Operational and priority distinctions]\label{rem:boundaryscope68}
Input erasure is decisive in \eqref{eq:productreplacer68}. No analogous
product-state reduction has been proved here for an arbitrary disturbing
instrument; coherent accumulation for unitary channels remains governed
by Proposition~\ref{prop:boundaryphase67}. The present theorems classify
reusable description lengths in the stated preparation families, not
general simulator workspace. Constants remain dimension dependent.

Quantum population compression starts with physical copies of an
unknown state and compresses their quantum information. Yang, Bai,
Chiribella and Hayashi \cite{YBCH68} obtain a half-parameter-count
logarithmic memory law in that different task and show a limitation
on purely classical compression of noncommuting families. Here the
encoder is given a classical matrix description, and each codeword
specifies a state or preparation instrument. It need not infer that
word from unknown copies. The common half-dimension scale, purification
inequalities and finite-dimensional covering methods are classical
antecedents, not claimed inventions. The contribution established here
is a uniform two-parameter legal-description law with a rational,
zero-outcome-preserving, rank-nonincreasing encoder and matching lower
bounds throughout the declared boundary families.
\end{remark}

\section{Retraction of covering centres for preparation targets}\label{sec:centres70}
The centres in a covering problem need not belong to the family being
covered. For preparation targets this distinction can be removed exactly,
without a small-error or rank assumption. Throughout this section a centre
is one memoryless instrument, reused at every call; it is not a device with
an additional persistent environment.

\begin{theorem}[Preparation-centre retraction]\label{thm:retraction70}
Let $\mathcal C=(\mathcal C_y)_{y=1}^m$ be any instrument from $M_d$ to
$M_n$, and set
\begin{equation}\label{eq:retraction70}
 \tau_y=\mathcal C_y(I_d/d)=d^{-1}\operatorname{Tr}_{\rm in}J_y(\mathcal C),
 \qquad \mathfrak R\mathcal C=\mathcal R_\tau.
\end{equation}
Then $\mathfrak R$ maps legal instruments to preparation instruments,
fixes every preparation instrument, and preserves rational Choi data.
For every $N\ge1$ and every $\sigma\in\mathcal S_{m,n}$,
\begin{equation}\label{eq:centredominance70}
 d_N(\mathcal R_\sigma,\mathfrak R\mathcal C)
                  \le d_N(\mathcal R_\sigma,\mathcal C).
\end{equation}
Here $d_N$ is the unhalved trace distance optimized over the common
adaptive testers of Lemma~\ref{lem:productreplacer68}. Consequently, for
any subset of preparation targets and any radius $\delta>0$, the minimum
covering cardinality is unchanged when arbitrary legal memoryless centres
are restricted to preparation centres. The assertion also holds when
centres are required to have rational Choi entries.
\end{theorem}
\begin{proof}
Complete positivity gives $\tau_y\ge0$, and trace preservation gives
$\sum_y\tr\tau_y=1$. The formula is rational and is the identity on the
preparation family, so $\mathfrak R^2=\mathfrak R$. Feed the state $I_d/d$
into each of $N$ calls to $\mathcal C$, use a fresh input at the next call,
and retain all outputs and outcome labels. The resulting state is
$\boldsymbol\tau^{\otimes N}$, since the centre is memoryless. The same
tester applied to $\mathcal R_\sigma$ returns
$\boldsymbol\sigma^{\otimes N}$. Therefore
\[
 \|\boldsymbol\sigma^{\otimes N}-\boldsymbol\tau^{\otimes N}\|_1
                        \le d_N(\mathcal R_\sigma,\mathcal C).
\]
Lemma~\ref{lem:productreplacer68} identifies the left side with the left
side of \eqref{eq:centredominance70}. Retracting every centre in any cover
preserves coverage and cannot increase cardinality. The reverse inequality
follows because preparation centres are among the legal centres. The same
argument preserves rationality.
\end{proof}

The retraction does not preserve a target rank bound: for example the
identity channel retracts to the maximally mixed preparation. Thus the
centre restriction in Theorem~\ref{thm:retraction70} is to the full
preparation family, not necessarily to the target rank stratum. The
rank-constrained rational upper centres of
Theorem~\ref{thm:factorcode68} remain a separate, stronger construction.
The theorem concerns the geometry of mathematical descriptions. Its
maximally mixed test is a quantum experiment, not a finite-classical-message
simulation of unknown inputs.

The reduction mechanism for input-erasing channels has a direct antecedent
in Cooney--Mosonyi--Wilde \cite{CMW70}. Their discussion of two replacer
channels reduces channel divergences to state divergences; their principal
results concern asymptotic discrimination and strong-converse exponents
for a channel versus a replacer. Lemma~\ref{lem:productreplacer68} uses the
same input-erasure mechanism in an exact finite-use trace-norm identity,
whereas Theorem~\ref{thm:retraction70} concerns the allowed centres of an
entire covering problem. No exact nonadaptive optimality for an arbitrary
non-replacer channel against a replacer is inferred from that identity.
The processor in the product identity is quantum, and input erasure,
rather than entanglement breaking alone, is its essential hypothesis.

\section{Coherent and preparation directions at a singular boundary}\label{sec:coherent70}
The preparation theorem has a square-root horizon scale. Reversible
quantum information can instead accumulate a coherent discrepancy at a
linear scale. We now exhibit both scales in a single non-input-erasing
family, and determine their combined description cost.

Fix $d\ge2$, $n,m\ge1$ and public rank bounds $0\le r_y\le n$, not all
zero. For $U\in U(d)$ and $\sigma\in\mathcal S_{\boldsymbol r}$ define
an instrument from $M_d$ to $M_{dn}$ by
\begin{equation}\label{eq:coherentfamily70}
 \mathcal F_{U,\sigma,y}(X)=UXU^*\otimes\sigma_y.
\end{equation}
Its outcome register is retained. Global phase is immaterial, so the
unitary parameter lies in $PU(d)$. Put
\[
 u=d^2-1,\qquad v=\sum_y r_y(2n-r_y)-1.
\]
The Choi matrix of each outcome is, up to a fixed tensor permutation,
$|U\rangle\!\rangle\langle\!\langle U|\otimes\sigma_y$. Its rank is
$\operatorname{rank}\sigma_y$, not $d\operatorname{rank}\sigma_y$ as in
the input-erasing family. These instruments transmit the data system
unitarily, including its entanglement with a reference. Their preparation
supports may change and may lie on any of the permitted singular strata.
Nevertheless, the outcome probabilities do not depend on the input. This
is a specified boundary family, not the whole body of disturbing
instruments.

Let $\mathcal C_N(\delta;\boldsymbol r)$ denote its covering number in
$d_N$, with centres allowed to be any legal memoryless instrument on the
same input, output and outcome spaces. Let $\mathcal C_N^{\Q}$ impose the
additional condition that the centres have Gaussian-rational Choi entries.
A target is supplied as $U$ and its preparation blocks; no unknown-device
learning interface is used.

\begin{theorem}[Joint coding at a coherent--preparation boundary]\label{thm:coherententropy70}
There are $0<c\le C<\infty$, depending only on the fixed dimensions and
rank bounds, such that, for every $N\ge1$ and $0<\delta\le1/32$,
\begin{equation}\label{eq:coherententropy70}
 c\left(\frac N\delta\right)^u
  \left(\frac{\sqrt N}\delta\right)^v
 \le\mathcal C_N(\delta;\boldsymbol r)
 \le\mathcal C_N^{\Q}(\delta;\boldsymbol r)
 \le C\left(\frac N\delta\right)^u
  \left(\frac{\sqrt N}\delta\right)^v.
\end{equation}
The metric is the unhalved final-state trace norm, maximized over common
adaptive quantum testers using the same memoryless instrument at most
$N$ times, with arbitrary finite references, quantum memory, feedback and
bounded public stopping. In particular, the optimal fixed-length reusable
description has length
\begin{equation}\label{eq:coherentbits70}
 (u+v/2)\log_2N+(u+v)\log_2(1/\delta)+O_{d,n,m,\boldsymbol r}(1).
\end{equation}
The upper centres can be chosen within the family
\eqref{eq:coherentfamily70}, with rational unitary matrices and rational
preparation blocks. They do not increase a target block rank and preserve
its zero outcomes. A finite exact encoder exists for supplied rational
unitaries and rational preparation blocks. Description length is not its
workspace or running time, and decoding describes a quantum instrument,
not a classical physical processor for unknown entangled inputs.
\end{theorem}

\subsection{The two metric scales}
For $U,V\in U(d)$ write
\begin{equation}\label{eq:chordunitary70}
 \chi(U,V)=\left(1-\frac{|\tr(U^*V)|^2}{d^2}\right)^{1/2},\qquad
 \Delta_N(\sigma,\tau)=
       \|\boldsymbol\sigma^{\otimes N}-\boldsymbol\tau^{\otimes N}\|_1.
\end{equation}
The first expression is the chordal distance between normalized Choi rays;
it is a metric after quotienting global phase.

\begin{lemma}[Anisotropic adaptive comparison]\label{lem:coherentmetric70}
For the instruments \eqref{eq:coherentfamily70},
\begin{align}
 \max\{\min\{N\chi(U,V)/\sqrt2,1\},\Delta_N(\sigma,\tau)\}
  &\le d_N(\mathcal F_{U,\sigma},\mathcal F_{V,\tau}),\label{eq:coherentlower70}\\
 d_N(\mathcal F_{U,\sigma},\mathcal F_{V,\tau})
  &\le\min\{2,2dN\chi(U,V)+\Delta_N(\sigma,\tau)\}.
                                                        \label{eq:coherentupper70}
\end{align}
All lower testers are permitted to depend on the pair being compared.
\end{lemma}
\begin{proof}
Discarding the data-system outputs gives the product-state lower bound.
For the unitary bound, discard the preparation systems and outcome labels,
and apply the same known operation $U^*$ to the data output after each
call. The two resulting data operations are $I$ and $W=U^*V$.
If $e^{i\phi_j}$ are the eigenvalues of $W$, then
\[
 \chi(U,V)^2=\frac2{d^2}\sum_{i,j}
                  \sin^2((\phi_i-\phi_j)/2).
\]
There is an eigenvalue pair whose principal half-angle
$\theta\in[0,\pi/2]$ satisfies $\sin\theta\ge\chi(U,V)/\sqrt2$.
Prepare the equal superposition of its orthogonal eigenvectors. After
$k$ calls its two possible pure states have trace distance
$2|\sin(k\theta)|$. If $N\theta\le\pi/2$, take $k=N$ and use
$\sin x\ge2x/\pi$ on $[0,\pi/2]$. The result is at least
$N\sin\theta$. If $N\theta>\pi/2$, take
$k=\lceil\pi/(4\theta)\rceil\le N$. Then
$\pi/4\le k\theta\le3\pi/4$, giving distance at least $\sqrt2$.
The case $\theta=0$ is immediate. This proves the first minimum in
\eqref{eq:coherentlower70}; the two lower witnesses need not coincide.

For the upper bound first change $U$ to $V$, keeping $\sigma$ fixed.
The unnormalized Choi difference of the unitary channels has trace norm
$2d\chi(U,V)$ and bounds their diamond distance. Appending a fixed
preparation state does not increase that norm. A slot-by-slot hybrid and
trace contraction for the common tester therefore cost at most
$2dN\chi(U,V)$. Next hold $V$ fixed and change $\sigma$ to $\tau$.
All calls to $V$ and all tester operations form one common quantum
processor consuming $N$ preparation states; unused copies are discarded
at a public stop. This costs at most $\Delta_N(\sigma,\tau)$, not the
sum of $N$ one-copy trace distances. Triangle inequality gives
\eqref{eq:coherentupper70}.
\end{proof}

Finite-use amplification of unitary differences is classical; see
Ac\'in \cite{Acin70} and Duan--Feng--Ying \cite{DFY70}. The elementary
estimate above is included to specify its uniform constant, its
pair-dependent controls and its compatibility with bounded stopping.
Neither the existence of unitary discrimination nor the real dimension
$d^2-1$ is claimed as a new ingredient.

\subsection{A finite rational atlas for unitary channels}
For $\ell\ge1$, a signed stereographic chart on $S^\ell$ has the form
\begin{equation}\label{eq:sphereatlas70}
 S_{j,\xi}(t)_j=\xi\frac{1-\|t\|^2}{1+\|t\|^2},\qquad
 S_{j,\xi}(t)_i=\xi\frac{2t_i}{1+\|t\|^2}\quad(i\ne j),
 \qquad \xi\in\{-1,1\}.
\end{equation}
For a target $q\in S^\ell$, choose $j$ of largest absolute coordinate,
$\xi=\operatorname{sgn}(q_j)$ and
$t_i=\xi q_i/(1+|q_j|)$. Then $|t_i|\le1$ and
$S_{j,\xi}(t)=q$. Direct expansion gives
\begin{equation}\label{eq:stereodistance70}
 \|S_{j,\xi}(t)-S_{j,\xi}(t')\|
 =\frac{2\|t-t'\|}{\sqrt{(1+\|t\|^2)(1+\|t'\|^2)}}.
\end{equation}
Rounding each coordinate towards zero at denominator $B$ thus gives a
rational unit vector at distance at most $2\sqrt\ell/B$. Every decoded
chart word is on the sphere, regardless of whether it is the canonical
encoding of a target.

\begin{lemma}[Rational unitary code]\label{lem:unitaryatlas70}
Put $L=d(d-1)/2$ and $K_d=3L+2(d-1)$. There is an explicitly enumerable
set $\mathcal U_B\subset U(d,\Q(i))$ satisfying
\begin{equation}\label{eq:unitarycode70}
 |\mathcal U_B|\le6^L4^{d-1}(2B+1)^{d^2-1},\qquad
 \sup_U\min_{V\in\mathcal U_B}\chi(U,V)\le K_d/B.
\end{equation}
For rational $U$, membership in any prescribed rational squared-error
ball $\chi(U,V)^2\le b$ is decidable by rational arithmetic alone.
\end{lemma}
\begin{proof}
The usual complex Givens elimination, in a fixed order of index pairs,
expresses every unitary as a product of $L$ embedded matrices
\[
 \begin{pmatrix}c&-\overline s\\s&c\end{pmatrix},
       \qquad c\in\R,\quad c^2+|s|^2=1,
\]
and a diagonal unitary. For completeness, a column pair $(a,b)$ is sent
to one nonzero coordinate using $c=|a|/(|a|^2+|b|^2)^{1/2}$ and
$s=cb/a$ when $a\ne0$; the cases $a=0$ or $(a,b)=0$ are obtained by
$c=0$ or by the identity. Applying the adjoint rotation eliminates the
second coordinate. Successive eliminations give an upper triangular
unitary, necessarily diagonal. Removing the phase of its first diagonal
entry leaves $d-1$ free phases. All factors are taken in the corresponding
fixed reversed order when reconstructing $U$.

Each Givens matrix has an $S^2$ parameter $(c,\Re s,\Im s)$; its
operator-norm change equals the Euclidean change of that vector. Each
remaining phase has an $S^1$ parameter. Apply
\eqref{eq:sphereatlas70} at grid $B$ to these $L+d-1$ factors.
There are $6$ chart headers for each $S^2$ and $4$ for each $S^1$,
with $2L+d-1=d^2-1$ grid digits altogether. Telescoping a product of
unitaries gives, up to a global phase, operator error at most
\[
 (2\sqrt2 L+2(d-1))/B\le K_d/B.
\]
For any fixed phase $z$ of modulus one,
$\chi(U,V)^2\le d^{-1}\|U-zV\|_F^2\le\|U-zV\|_{\rm op}^2$.
This proves \eqref{eq:unitarycode70}. The formula
\eqref{eq:chordunitary70} has rational square for rational $U,V$, proving
the last assertion. Enumeration of these finite words followed by the
squared-distance test is an exact encoder whenever the radius is at least
$K_d/B$. This does not assert a fast nearest-centre algorithm.
\end{proof}

\subsection{The joint covering theorem}
\begin{proof}[Proof of Theorem~\ref{thm:coherententropy70}]
For the upper bound use
\[
 B_U=\left\lceil4dK_dN/\delta\right\rceil,
 \qquad B_P=\max\{1,\lceil8\sqrt{Nv}/\delta\rceil\}.
\]
Lemma~\ref{lem:unitaryatlas70} makes the first term in
\eqref{eq:coherentupper70} at most $\delta/2$.
The factor code in Theorem~\ref{thm:factorcode68} makes the second term
at most $\delta/2$. Its finite pivot headers and anchor code give at most
$C_{n,m,\boldsymbol r}(2B_P+1)^v$ legal images. If $v=0$ the preparation
family is a singleton and no such digits are needed. Multiplication of
the two cardinalities proves the upper bound and all legality, rank and
zero-outcome assertions. For rational data, search the finite unitary
list using the rational test
$\chi(U,V)^2\le\delta^2/(16d^2N^2)$; the preceding estimate guarantees
a result. The preparation encoder is already exact. Decoded Choi entries
are rational products of the two rational outputs. The finite search and
the expanded matrices are not charged to the reusable payload.

For the converse, a neighborhood of the identity in $PU(d)$ contains a
quantitatively bi-Lipschitz $u$-dimensional Euclidean patch for the metric
$\chi$. Here is a direct verification. Use traceless Hermitian coordinates
$H$ near zero and the ray of $e^{iH}$. The differential of its normalized
Choi projection at zero vanishes only when $H$ is scalar, hence only when
$H=0$. On a sufficiently small convex coordinate ball its differential
is uniformly closer than half its least singular value to that at zero.
Integration along a line segment gives both Lipschitz inequalities.
Since the projection Frobenius distance is $\sqrt2\chi$, the asserted
metric comparison follows. Volumetric packing therefore gives at least
$c_d(N/\delta)^u$ unitaries separated by more than $4\delta/N$, with
a smaller fixed constant covering the coarse-scale cases as well.

The maximal-rank preparation patch used in
Theorem~\ref{thm:rankentropy68} has real dimension $v$ and is
bi-Lipschitz in tuple Frobenius norm. Pack it at separation greater than
$32\delta/\sqrt N$, obtaining at least
$c_{n,m,\boldsymbol r}(\sqrt N/\delta)^v$ points. The local binary
bound in Corollary~\ref{cor:localchoi68}, applied to the preparation
state (input dimension one), separates their $N$-copy states by more
than $2\delta$. Indeed the bound at that scale is
$(3/16)\min\{16\delta,1\}$, which is $3\delta$ before saturation
and at least $3/16>2\delta$ afterwards. For $v=0$ use its single point.

Take the Cartesian product of these packings. Two elements with different
unitaries are separated by \eqref{eq:coherentlower70}, since
$\min\{2\sqrt2\delta,1\}>2\delta$. Two elements with the same unitary
and different preparation blocks are separated by the other lower witness.
Thus a ball of radius $\delta$ about any legal centre can contain at most
one product point. This proves the lower bound even for centres outside
the family. Taking base-two logarithms and rounding a fixed-length word
length upward changes the answer by at most a fixed constant, proving
\eqref{eq:coherentbits70}.
\end{proof}

The coefficient $u+v/2$ reflects two different reuse scales, rather than
half of the total parameter dimension $u+v$. No spectral-gap assumption
or positive eigenvalue margin is used. The constants are nevertheless
fixed-dimensional. Nor is this a classification of every tangent direction
of an arbitrary instrument: coupling an input-dependent measurement to
its disturbance is not part of \eqref{eq:coherentfamily70}.

\subsection{An explicit qubit codec}
For $q=(q_0,q_1,q_2,q_3)\in S^3$ put
\begin{equation}\label{eq:quaternion70}
 U(q)=\begin{pmatrix}q_0+iq_3&q_2+iq_1\\-q_2+iq_1&q_0-iq_3\end{pmatrix}.
\end{equation}
The points $q$ and $-q$ describe the same channel. Choose a largest
absolute coordinate, flip that common sign to make it positive, and use
\eqref{eq:sphereatlas70} with $\ell=3$ and $\xi=1$. Its header has
four possibilities, and its body has exactly three signed digits.

\begin{corollary}[Exact rational coherent--preparation description]\label{cor:qubitcodec70}
For a supplied rational unit quaternion and rational preparation blocks,
the preceding three-digit chart, at $B_U=\lceil16N/\delta\rceil$, combined
with the preparation codec at tolerance $\delta/2$, gives a legal rational
instrument within $d_N$ error $\delta$. Its fixed-length description has
\[
 (3+v/2)\log_2N+(3+v)\log_2(1/\delta)+O_{n,m,\boldsymbol r}(1)
\]
bits, with no increase of a block rank and exact preservation of zero
outcomes. It uses only integer and rational arithmetic, including the
integer square roots in the preparation code. Its output is a mathematical
instrument description, not physical classical simulation.
\end{corollary}
\begin{proof}
For quaternion matrices,
$\|U(q)-U(q')\|_{\rm op}=\|q-q'\|$. The unitary diamond error is
at most twice this quantity. Thus the chart's coherent contribution
is at most $4\sqrt3 N/B_U<\delta/2$. Its squared upper certificate is
$48N^2/B_U^2$. The preparation contribution is at most $\delta/2$.
Triangle inequality proves the claim; one may certify its square by
$2(e_U^2+e_P^2)$, not by $e_U^2+e_P^2$. The chart returns an exactly
unit quaternion, and the preparation decoder returns positive blocks
of total trace one. Formula~\eqref{eq:coherentfamily70} therefore gives
exact complete positivity and trace preservation, including on singular
inputs. The chart has at most $4(2B_U+1)^3$ words, which gives the
stated length after adding the inherited preparation payload.
\end{proof}
The reference encoder replays both deterministic encoders when verifying
a supplied target. Bare decoding checks only syntax, rational legality and
the declared construction budget. It cannot certify proximity to an
unspecified target. In particular, reference stability concerns the quantum
instrument defined by the decoded Choi matrices, not access to a classical
processor with equivalent action on an unknown reference-entangled state.

\section{Coverings throughout the submaximal error range}\label{sec:fullerror71}
A pairwise separation argument at radius $\delta$ requires distances
larger than $2\delta$. Since our adaptive distance is an unhalved trace
norm, this argument cannot by itself reach $\delta\ge1$. We instead
bound the number of targets that one centre can cover. The same
measurement, rather than a different binary test for each pair, will be
used in this part of the converse. Dimensions are fixed throughout. For a target family $\mathcal F$, write
$\mathcal M_N(\delta;\mathcal F)$ for its least covering cardinality
in $d_N$ with arbitrary legal memoryless centres, and add a superscript
$\Q$ when the decoded Choi blocks must be rational.

\begin{lemma}[Multiplicity of a covering centre]\label{lem:multiplicity71}
Let $J_1,\ldots,J_K$ be targets and let one common tester with at most
$N$ calls have disjoint classical events $A_1,\ldots,A_K$ such that
$\Pr_{J_i}(A_i)\ge1-\eta$. If $0\le\eta<1$ and
$0\le\delta<2(1-\eta)$, every covering of these targets by $M$ legal
centres at adaptive radius $\delta$ satisfies
\begin{equation}\label{eq:multiplicity71}
                M\ge (1-\eta-\delta/2)K.
\end{equation}
A centre need not belong to the target family. Indeed the conclusion
still holds for a centre specified by a legal $N$-slot process with
private memory, provided its distance to each assigned target controls
this same tester.
\end{lemma}
\begin{proof}
Let $Q$ be the tester's output distribution at a centre. For each target
assigned to that centre, trace contraction under the terminal
measurement gives
$Q(A_i)\ge1-\eta-\delta/2$. The division by two is required because
classical total variation is half the unhalved trace norm. Disjointness
gives at most $(1-\eta-\delta/2)^{-1}$ assigned targets. Sum this
bound over centres. No product structure, unbiasedness or independence
is required of $Q$.
\end{proof}

\begin{lemma}[One common finite-dimensional estimator]\label{lem:tomography71}
Put $b=m(dn)^2$. On $\mathcal C_{d,n,m}$ there is, for every $N\ge1$,
one nonadaptive tester and a Hermitian tuple estimator $\widehat J$
satisfying
\begin{equation}\label{eq:tomography71}
 \sup_{J\in\mathcal C_{d,n,m}}
      \E_J |\widehat J-J|_2^2\le \frac{C_0}{N},
                  \qquad C_0=2d^2b^2.
\end{equation}
The estimator need not be positive or trace preserving. The tester
and estimator are independent of the unknown target and of any later
choice of a packing set.
\end{lemma}
\begin{proof}
One normalized maximally entangled input produces the flagged Choi
state $\Omega_J=d^{-1}\bigoplus_y J_y$. Choose a Hilbert--Schmidt
orthonormal Hermitian basis $H_1,\ldots,H_b$ of the block-diagonal
space. Every $H_j$ has operator norm at most one. Thus
$(I+H_j)/2,(I-H_j)/2$ is a binary POVM, and its signed outcome has
mean $\operatorname{tr}(H_j\Omega_J)$ and variance at most one.
For $N\ge b$, use $k=\lfloor N/b\rfloor$ independent Choi inputs
for each of the $b$ measurements. There are at most $N$ calls. Let
$\widehat a_j$ be the corresponding empirical signed mean and set
$\widehat\Omega=\sum_j\widehat a_jH_j$. Orthonormality gives
\[
 \E_J\|\widehat\Omega-\Omega_J\|_F^2
   =\sum_j\operatorname{Var}(\widehat a_j)
   \le b/k\le2b^2/N.
\]
Set $\widehat J=d\widehat\Omega$, interpreted as a tuple. For $N<b$
use $\widehat J=0$ and make no calls. Then
$|J|_2^2=d^2\operatorname{tr}\Omega_J^2\le d^2\le C_0/N$.
This proves every asserted range, without a rank or eigenvalue-margin
condition. The binary POVMs are block diagonal in the public outcome,
so coherent access to different classical flags is not needed.
\end{proof}
This elementary informationally complete estimation argument is a
standard finite-dimensional tomography construction; compare the process
estimation setting of \cite{PLS66}. Here it is used only to build a common
lower-bound witness. It does not constitute a new optimal-learning
claim, and the description encoder below is still supplied its target.

\begin{theorem}[A full-range covering lower bound]\label{thm:fullpacking71}
Let $\mathcal F\subseteq\mathcal C_{d,n,m}$ contain the image of a
nonempty $a$-dimensional Euclidean ball under a bi-Lipschitz map for
tuple Frobenius distance, with $a>0$. For every fixed $0<\delta_*<2$
there is $c>0$, depending on this patch, the fixed dimensions and
$\delta_*$, such that
\begin{equation}\label{eq:fullpacking71}
       \mathcal M_N(\delta;\mathcal F)\ge cN^{a/2}
             \quad(N\ge1,\ 0<\delta\le\delta_*).
\end{equation}
Centres may be arbitrary legal memoryless instruments. The same lower
bound holds under the relaxation to legal $N$-slot centres described
in Lemma~\ref{lem:multiplicity71}.
\end{theorem}
\begin{proof}
Set $\eta=(2-\delta_*)/4>0$ and
$r_N=(C_0/(\eta N))^{1/2}$. The previous lemma and Markov's inequality
give
$\Pr_J\{|\widehat J-J|_2\le r_N\}\ge1-\eta$ for every target.
A volumetric packing of the fixed patch supplies at least
$c_1r_N^{-a}$ points at pairwise tuple distance greater than $2r_N$.
The constant can be decreased to cover coarse scales: a singleton
packing always exists. The estimator events consisting of the closed
$r_N$-balls at these points are disjoint. Lemma~\ref{lem:multiplicity71}
applies because
$1-\eta-\delta/2\ge1-\eta-\delta_*/2=\eta$.
It gives $\mathcal M_N\ge\eta c_1r_N^{-a}$, as required. In
particular one centre may cover several packing points when the error
is large; the proof bounds that multiplicity rather than asserting it
is one.
\end{proof}

\begin{corollary}[Full-range interior and preparation laws]\label{cor:fullrange71}
Fix $0<\delta_*<2$. In Theorem~\ref{thm:jointcode68}, with fixed
$d,n,m,a$ and $s=d^2(mn^2-1)>0$, both real and rational legal-centre
covering numbers satisfy
\begin{equation}\label{eq:interiorfull71}
 \log_2\mathcal M_N(\delta;a)
 =\log_2\mathcal M_N^{\Q}(\delta;a)+O(1)
 =\frac{s}{2}\log_2N+s\log_2(1/\delta)+O(1)
\end{equation}
uniformly for $N\ge1$ and $0<\delta\le\delta_*$. The constants may
depend on $\delta_*$ and the positive margin $a$.
For the preparation family in Theorem~\ref{thm:rankentropy68}, the
same extension holds with $s$ replaced by
$v=\sum_y r_y(2n-r_y)-1$, with fixed dimensions and public output-rank
bounds. Its constants require no positive eigenvalue margin. If $v=0$
the family is a singleton and needs no payload.
\end{corollary}
\begin{proof}
For $0<\delta\le1/32$ these are the inherited joint laws. The interior
ball and maximal-rank preparation chart used in their proofs are
bi-Lipschitz patches of dimensions $s$ and $v$, respectively. Theorem
\ref{thm:fullpacking71} therefore supplies the lower term
$(s/2)\log_2N+O(1)$, or $(v/2)\log_2N+O(1)$, on the remaining
range $1/32<\delta\le\delta_*$. On that range
$\log_2(1/\delta)$ is bounded above and below by constants depending
only on $\delta_*$. For the upper bound use the old code at accuracy
$1/32$, which is also accurate to $\delta$. This proves a uniform
additive remainder for both kinds of centres. If $\delta_*\le1/32$
there is no remaining range.
\end{proof}
The cap $\delta_*<2$ is substantive: at unhalved error two every
nonempty family is covered by a single legal centre. Nothing above is
uniform as $\delta_*\uparrow2$, or along a varying error sequence
converging to two. The exact zero-error endpoint is also separate from
the positive-error coding law. The older coherent/preparation theorem
retains its own $\delta\le1/32$ range; the present common tomography
argument gives the half-dimension scale, not its coherent $N^{-1}$ scale.

\section{Input-dependent boundary instruments}\label{sec:cqboundary71}
Fix an orthonormal basis $e_1,\ldots,e_d$ of the input and the associated
complete dephasing map
$\mathcal D(X)=\sum_x\langle e_x,Xe_x\rangle|e_x\rangle\langle e_x|$.
For positive semidefinite $n\times n$ matrices $\sigma_{xy}$ satisfying
$\sum_y\operatorname{tr}\sigma_{xy}=1$ for every $x$, define
\begin{equation}\label{eq:cqfamily71}
 \Phi_{\boldsymbol\sigma,y}(X)
       =\sum_{x=1}^d\langle e_x,Xe_x\rangle\sigma_{xy},
       \qquad 1\le y\le m.
\end{equation}
These instruments measure a fixed input basis and then prepare an
outcome-dependent state. Their public outcome probabilities
$\sum_x\langle e_x,Xe_x\rangle\operatorname{tr}\sigma_{xy}$ may
depend on the input. Measurement destroys input coherences; the family
is not in general input erasing, and it is not the coherent tensor
family of Theorem~\ref{thm:coherententropy70}.

The Choi blocks and their ranks are
\begin{equation}\label{eq:cqchoi71}
 J_y=\bigoplus_{x=1}^d\sigma_{xy},\qquad
 \operatorname{rank}J_y=\sum_x\operatorname{rank}\sigma_{xy}.
\end{equation}
In particular the marginal constraint is exactly one unit-trace
condition per input row. Fix public integers $0\le r_{xy}\le n$,
with $\sum_y r_{xy}>0$ in every row, and let $\mathcal F_{\boldsymbol r}$
be the family with $\operatorname{rank}\sigma_{xy}\le r_{xy}$. Put
\begin{equation}\label{eq:cqdim71}
 v_x=\sum_y r_{xy}(2n-r_{xy})-1,\qquad V=\sum_x v_x.
\end{equation}
The ranks here are conditional output-state ranks, not independent
rank bounds on arbitrary Choi blocks. The input basis, dimensions and
rank bounds are fixed, not learned from an unknown device.

\begin{theorem}[Rank-boundary coding with input-dependent outcomes]\label{thm:cqentropy71}
Fix the preceding data and $0<\delta_*<2$. If $V>0$, the least covering
cardinality by arbitrary legal memoryless centres and the least
cardinality with rational decoded Choi blocks both obey
\begin{equation}\label{eq:cqentropy71}
 \log_2\mathcal M_N(\delta;\mathcal F_{\boldsymbol r})
   =\frac V2\log_2N+V\log_2(1/\delta)+O(1)
\end{equation}
uniformly for $N\ge1$ and $0<\delta\le\delta_*$. The error is the
unhalved final-state trace distance over arbitrary common quantum
testers with finite references, quantum memory, feedback and public
stopping bounded by $N$. A rational upper code has decoded centres in
\eqref{eq:cqfamily71}; it never increases any conditional block rank
and preserves every zero block exactly. Its encoder on rational target
data uses integer and rational arithmetic only. If $V=0$ the family
is a known singleton.

The remainder depends on the fixed dimensions, rank bounds and
$\delta_*$, but not on $N$, $\delta$ or a smallest positive output
eigenvalue. Payload length is distinct from encoder workspace and
expanded decoder storage. The decoded object is a mathematical
quantum instrument, not a physical classical simulator on an unknown
reference-entangled input.
\end{theorem}

\subsection{Centres and the adaptive metric}
Write $\Omega_x=\bigoplus_y\sigma_{xy}$ for the flagged state of one
input row. A useful general observation first identifies the allowed
centres.

\begin{proposition}[Idempotent input retraction]\label{prop:inputretraction71}
Let $\Pi$ be a fixed completely positive trace-preserving idempotent
input map and let every target satisfy $\Phi\Pi=\Phi$. Then
\begin{equation}\label{eq:inputretraction71}
             d_N(\Phi,K\Pi)\le d_N(\Phi,K)
\end{equation}
for every legal memoryless centre $K$ and every $N$. Thus covering
numbers with all legal centres equal those with centres fixed by
right composition with $\Pi$, at every radius.
For $\Pi=\mathcal D$, the retracted centre has rows
$\tau_{xy}=K_y(|e_x\rangle\langle e_x|)$ and is of the form
\eqref{eq:cqfamily71}. In the prescribed coordinate basis this operation
preserves rationality. It need not preserve the target's rank bounds.
\end{proposition}
\begin{proof}
For any tester of $\Phi$ against $K\Pi$, insert $\Pi$ on its input
immediately before each use of the unknown device. This is a legal
common tester of $\Phi$ against $K$: in the first hypothesis
$\Phi\Pi=\Phi$, and in the second it implements $K\Pi$. References,
feedback and stopping are unaffected. Taking the supremum proves
\eqref{eq:inputretraction71}. Idempotence says that $K\Pi$ is fixed
by right composition with $\Pi$. Map every covering centre this way;
cardinality cannot increase. The reverse covering inequality follows
from inclusion of allowed centre classes. For dephasing, the displayed
rows are positive and have unit total trace, since $K$ is an instrument.
They are diagonal input blocks of the original Choi matrices.
\end{proof}
This is a data-processing retraction, not a rank-preserving projection.
It generalizes the preparation-centre observation of
Theorem~\ref{thm:retraction70}, whose fixed input map discards its input
and supplies one chosen state. In the converse of
Theorem~\ref{thm:cqentropy71}, no rank promise is imposed on centres.

\begin{lemma}[Row witnesses and a common programme]\label{lem:cqmetric71}
For two instruments of \eqref{eq:cqfamily71},
\begin{equation}\label{eq:cqsandwich71}
 \max_x\|\Omega_x^{\otimes N}-\Lambda_x^{\otimes N}\|_1
 \le d_N(\Phi_{\boldsymbol\sigma},\Phi_{\boldsymbol\tau})
 \le\left\|\bigotimes_x\Omega_x^{\otimes N}
               -\bigotimes_x\Lambda_x^{\otimes N}\right\|_1.
\end{equation}
If unit purifications $a_x,b_x$ of the corresponding row states are
chosen and $\eta_x=\|a_x-b_x\|$, then
\begin{equation}\label{eq:cqfactorbound71}
 d_N(\Phi_{\boldsymbol\sigma},\Phi_{\boldsymbol\tau})
                  \le2\sqrt{N\sum_x\eta_x^2}.
\end{equation}
For one use there is the exact identity
$d_1=\max_x\|\Omega_x-\Lambda_x\|_1$.
\end{lemma}
\begin{proof}
For the first inequality repeatedly feed $|e_x\rangle\langle e_x|$
and retain all flagged outputs. For the second, prepare, for each
of the $N$ slots, one copy of every row state. A common processor
measures the actual input in the prescribed basis and uses only the
programme state with that outcome $x$. The remaining row states in
that slot are discarded. If the input is entangled with the tester's
memory, its conditional memory operator is retained by this measurement;
appending the selected row state gives exactly
\eqref{eq:cqfamily71} on that joint input. The unobserved input outcome
is internal to the processor, not new feedback granted to the tester.
All later operations and the discard of unused slots at a stop form one
common CPTP processor. Trace contraction proves the upper bound.

Purify all these programme states by the chosen vectors. For unit
vectors, $1-|\langle a_x,b_x\rangle|^2\le\eta_x^2$. The elementary
product inequality $1-\prod t_j\le\sum(1-t_j)$ for $t_j\in[0,1]$
gives pure-programme trace distance squared at most
$4N\sum_x\eta_x^2$. Partial trace and the common processor give
\eqref{eq:cqfactorbound71}. Finally, for an arbitrary one-use input
with reference, the difference is
$\sum_x\xi_R^x\otimes(\Omega_x-\Lambda_x)$, where
$\xi_R^x\ge0$ and $\sum_x\operatorname{tr}\xi_R^x=1$.
The triangle inequality gives the one-use upper maximum, and a basis
input attains it.
\end{proof}
The product programme consumes up to $dN$ row states in the proof,
whereas the actual tester makes at most $N$ instrument calls. It is
only a comparison device, not an equality of these resource counts.
Neither inequality in \eqref{eq:cqsandwich71} is asserted to be an
equality for general $N$. In particular the absence of an asymptotic
adaptive advantage in a binary discrimination exponent does not imply
a finite-use metric equality; compare \cite[Sections~II and V]{SHW68}.

\subsection{Exact rational decoding and the covering proof}
\begin{lemma}[A common-grid row code]\label{lem:cqcodec71}
For rational target data, apply the singular factor encoder of
Theorem~\ref{thm:factorcode68} separately in every input row at a
common integer grid $B\ge1$. If $h_x$ is that row's actual real
factor-coordinate count, the decoded instrument satisfies
\begin{equation}\label{eq:cqcodeerror71}
 d_N(\Phi,\widetilde\Phi)^2
       \le\frac{16N\sum_x(h_x-1)}{B^2}
       \le\frac{16NV}{B^2}.
\end{equation}
There are at most $C_{d,n,m,\boldsymbol r}(2B+1)^V$ legal decoder
images. The conditional block ranks do not increase; consequently no
outcome Choi rank increases either. All zero conditional blocks remain
zero. The expanded common denominator is at most
$B^{2d}\prod_x h_x$ and thus has $O_{d,n,m}(1+\log B)$ bits.
\end{lemma}
\begin{proof}
The inherited anchor construction gives a unit factor purification in
row $x$ within distance $2\sqrt{h_x-1}/B$ of the target's chosen
unit factor vector. Apply \eqref{eq:cqfactorbound71}. Trace normalization
is performed separately in each row, so \eqref{eq:cqchoi71} is exactly
trace preserving. The inherited zero-pivot proof, rational squared
coordinates and integer-square-root ratio calculation apply row by row.
Their finite masks and anchor headers have fixed cardinality; the row
body has at most $v_x$ signed grid coordinates. Multiplying these
counts gives the asserted exponent $V$. Only valid decoded words count.
The block-rank and zero statements hold row by row and pass to the direct
sum in \eqref{eq:cqchoi71}. Row $x$ has normalization denominator
$T_x\le h_xB^2$. A common denominator is the least common multiple of
the $T_x$, bounded by their product. This proves the final assertion.
\end{proof}

\begin{proof}[Proof of Theorem~\ref{thm:cqentropy71}]
For the upper bound take
$B=\max\{1,\lceil4\sqrt{NV}/\delta\rceil\}$. The last lemma gives
error at most $\delta$ and at most $C(2B+1)^V$ centres, hence the
upper description law for $\delta\le\delta_*<2$. Real targets admit
the same factor rounding construction; the decoded matrices are still
rational. For supplied rational targets, no irrational factor is stored:
the exact inherited square-comparison algorithm determines the digits.
Thus the real-family rational-centre upper bound does not restrict the
targets to a countable subclass.

For the lower bound choose, in every row with $v_x>0$, the gauge-fixed
maximal-rank factor patch used in Theorem~\ref{thm:rankentropy68}.
Its inverse is the local leading-block Cholesky map. Rows with $v_x=0$
are fixed. A Cartesian product of sufficiently small closed coordinate
balls contains a $V$-dimensional ball. The map to the tuple
$(\sigma_{xy})$, and hence to $J$, is bi-Lipschitz; direct sum is an
isometry for tuple Frobenius norm. This supplies one fixed
$V$-dimensional patch with no claim of a global factor chart.
For $\delta\le1/32$ pack it at separation greater than
$32d\delta/\sqrt N$. Corollary~\ref{cor:localchoi68} gives distances
larger than $2\delta$, so even an unrestricted legal centre covers at
most one packing point. This gives
$c(\sqrt N/\delta)^V$ centres. For $1/32<\delta\le\delta_*$ use
Theorem~\ref{thm:fullpacking71} on the same patch and the boundedness
of $\log(1/\delta)$ on that interval. These lower bounds apply also
to rational centres, which form a smaller class. Taking logarithms,
with the usual ceiling for fixed-length words, completes the proof.
If $V=0$, each nonempty row consists of a specified one-dimensional
unit block; its location is fixed by the public rank bounds. Thus the
whole family is a known singleton.
\end{proof}
For example, with scalar output $n=1$ and all rank bounds equal to
one, \eqref{eq:cqentropy71} has $V=d(m-1)$: it is a boundary coding
law for the rows of a classical stochastic matrix. With one quantum
output and full row ranks, $V=d(n^2-1)$; with one pure output per row,
$V=d(2n-2)$. Mixed and pure conditional states need not commute.
These dimensions are standard rank-stratum dimensions; the theorem
identifies their joint coding scale under the declared adaptive metric.
The measured input basis is fixed. No law for varying measurement bases,
unknown instruments, or all irreversible quantum dynamics follows from
this result.

\section{Comparison and boundary geometry}
\subsection{Adaptive discrimination and description}
Salek--Hayashi--Winter \cite[Sections~II--V]{SHW68} study asymptotic
binary channel-discrimination exponents. Their results include equal
adaptive and nonadaptive exponents for classical--quantum channels,
and separations for some entanglement-breaking channels. Our object
is instead a covering number of a supplied family in the finite-use
metric $d_N$. Its lower packing may use a different binary test for
each pair; it does not supply a universal measurement which identifies
every codeword. The exact reduction in
Lemma~\ref{lem:productreplacer68} uses input erasure, not merely
entanglement breaking, and implies no general absence of an adaptive
advantage for entanglement-breaking channels.

Cooney--Mosonyi--Wilde \cite{CMW70} is a direct antecedent for the
input-erasure mechanism. Its two-replacer discussion and asymptotic
channel-versus-replacer exponents are distinct from an exact finite-use
covering problem. The retraction of Theorem~\ref{thm:retraction70} concerns
admissible centres, not an improvement to those discrimination exponents.
The coherent amplification used here likewise has the established unitary
discrimination antecedents \cite{Acin70,DFY70}.

The two-point programme method has a metrological antecedent in
\cite{DKG67}. Likewise the half-parameter scale has antecedents in
quantum population compression \cite{YBCH68}; there the input consists
of unknown physical copies and the memory interface may include
quantum storage. Here the encoder is given classical matrix data and
the decoder specifies an operation or state. These are different
coding tasks, even when a logarithmic coefficient has the same form.

Mele--Bittel \cite{MB67} and Oufkir--Girardi \cite{OG67} investigate
learning unknown channels in diamond distance. The present encoder
receives its target and asserts no query or sample bound. The
reference-sensitive physical classical-simulation problem of
Naik--Zartab--Gisin--Banik \cite{NZGB67} is different again: a rational
description with a diamond guarantee remains a description of a
quantum operation, not a finite-classical-message implementation
accepting an unknown entangled input. No bound here is asserted for
the hardware dimension of a universal programmable processor.

\subsection{Ranks, tangent directions and remaining joint regimes}
For a general instrument $J$, let $P_y$ project onto the support of
$J_y$. The one-sided tangent cone of the legal instrument body is
\[
 T_{\mathcal C}(J)=\overline{\{t(L-J):t\ge0,\ L\in\mathcal C\}}.
\]
Every element $H$ of this cone satisfies the marginal constraint and
$P_y^\perp H_yP_y^\perp\ge0$. These necessary conditions identify
where positivity becomes active; by themselves they assert no reuse
modulus. A boundary classification must examine the metric of
repeated use along support-changing directions and along reversible
coherent directions, rather than count affine coordinates alone.
Proposition~\ref{prop:boundaryphase67} shows that a coherent unitary
phase can have scale $N^{-1}$, whereas the preparation family in
Theorem~\ref{thm:rankentropy68} has an $N^{-1/2}$ coding scale even
at rank-deficient outputs. In that preparation family the Choi ranks
are $d\,\operatorname{rank}\sigma_y$, and the rank bounds in the
theorem are output ranks. They are not arbitrary Choi-rank strata of
all instruments.

Theorem~\ref{thm:coherententropy70} now gives a joint law in the
non-input-erasing tensor family $UXU^*\otimes\sigma_y$. It includes a
full projective unitary manifold and support-changing preparation blocks,
but no input-dependent measurement probabilities. Its mixed coefficient
$u+v/2$ is established by two explicit adaptive moduli and matching product
packings, not by a tangent-dimension count alone. These qualifications are
essential: it is not a classification of all coherent/noisy couplings.

Theorems~\ref{thm:jointcode68} and~\ref{thm:rankentropy68} settle the
joint $N,\delta$ dependence in their specified fixed-dimensional
families. They do not assert dimension-uniform constants, a uniform
limit as the positive margin tends to zero, or a full tangent-cone
classification for general disturbing instruments. The older
exponential-accuracy width result retains its multiplicative gap in
the separate label-width resource. None of these description laws
is a mutable-workspace converse for general variable-input
trajectory simulation.


\subsection{Input dependence and a common full-range witness}
The fixed-readout family of Theorem~\ref{thm:cqentropy71} is a
classical--quantum channel with a fixed quantum input measurement.
Its $N$-slot programme selects conditional row states and therefore
handles input-dependent probabilities. This is within the channel
families considered by Salek--Hayashi--Winter
\cite[Sections~II and V]{SHW68} for asymptotic binary discrimination.
Their exponent comparisons are direct antecedents, but are not a
finite-use family covering number or a rational rank-aware code.
We do not infer equality between the row-product upper bound and the
adaptive distance, nor a finite-use prohibition on adaptive advantage.
The outcome/state rows may be noncommuting and have zero blocks.
The measured basis remains fixed, so a varying measurement basis and
general coherent/noisy couplings are not parameterized by $V$.

The common-estimator argument in Section~\ref{sec:fullerror71} uses
standard informationally complete tomography \cite{PLS66} and an
elementary event-counting inequality. Its point is to control covering
multiplicity at error greater than one, where the earlier requirement
of pairwise distance greater than $2\delta$ is unavailable. This new
witness is deliberately common to all packing points; the local
small-error proof still uses pair-dependent binary tests. Neither
witness turns the supplied-description code into an unknown-instrument
learning result. The full-range extension concerns half-dimensional
coding scales and does not extend the coherent $N^{-1}$ scale by fiat.
The author-side comparisons here identify premises and conclusions;
they do not constitute an independent exhaustive priority review.

\begin{thebibliography}{99}
\bibitem{Watrous65} J. Watrous, \emph{The Theory of Quantum Information},
Cambridge University Press, 2018. Theorems~2.22 and~2.26
(Choi criteria), Corollary~3.40 (trace-norm contraction),
and Section~3.3, equation~(3.306) (hybrid comparison).
\bibitem{PLS66} T. Surawy-Stepney, J. Kahn, R. Kueng and M. Gu\c{t}\u{a},
Projected least-squares quantum process tomography,
\emph{Quantum} \textbf{6} (2022), 844;
arXiv:2107.01060v2, 18 October 2022.
\bibitem{Gutoski66} G. Gutoski,
On a measure of distance for quantum strategies,
\emph{J. Math. Phys.} \textbf{53} (2012), 032202;
arXiv:1008.4636v4, 14 March 2012.
\bibitem{NZGB67} S. G. Naik, M. Zartab, N. Gisin and M. Banik,
No-Go Theorem for Generic Simulation of Qubit Channels with Finite
Classical Resources, arXiv:2501.15807v2 (16 July 2025),
Theorems~2 and~5.
\bibitem{MB67} A. A. Mele and L. Bittel,
Optimal learning of quantum channels in diamond distance,
arXiv:2512.10214v3 (2026), Theorem~I.1 and Sections~III.1.1, IV.1.1.
\bibitem{OG67} A. Oufkir and F. Girardi,
Improved Lower Bounds for Learning Quantum Channels in Diamond
Distance, arXiv:2601.04180v4 (2026).
\bibitem{DKG67} R. Demkowicz-Dobrza\'nski, J. Ko\l ody\'nski
and M. Gu\c{t}\u{a}, The elusive Heisenberg limit in quantum-enhanced
metrology, \emph{Nature Communications} \textbf{3} (2012), 1063.
DOI: 10.1038/ncomms2067; arXiv:1201.3940v2, equations~(4)--(6)
and the classical-simulation argument in Methods.
\bibitem{YBCH68} Y.~Yang, G.~Bai, G.~Chiribella and M.~Hayashi,
Compression for quantum population coding,
\emph{IEEE Transactions on Information Theory} \textbf{64} (2018), no.~7,
doi:10.1109/TIT.2017.2788407;
arXiv:1701.03372v8.
\bibitem{SHW68} F.~Salek, M.~Hayashi and A.~Winter,
Usefulness of adaptive strategies in asymptotic quantum channel
 discrimination, \emph{Phys. Rev. A} \textbf{105} (2022), 022419;
arXiv:2011.06569v3 (17 March 2024).
\bibitem{CMW70} T.~Cooney, M.~Mosonyi and M.~M.~Wilde,
Strong converse exponents for a quantum channel discrimination problem
and quantum-feedback-assisted communication,
arXiv:1408.3373v3 (26 February 2016), p.~5 and Sections~4.1--4.2.
\bibitem{Acin70} A.~Ac\'in,
Statistical distinguishability between unitary operations,
\emph{Phys. Rev. Lett.} \textbf{87} (2001), 177901.
DOI: 10.1103/PhysRevLett.87.177901.
\bibitem{DFY70} R.~Duan, Y.~Feng and M.~Ying,
Entanglement is not necessary for perfect discrimination between
unitary operations, \emph{Phys. Rev. Lett.} \textbf{98} (2007), 100503.
DOI: 10.1103/PhysRevLett.98.100503.
\end{thebibliography}

\end{document}
